\documentclass[11pt]{article}
\usepackage[a4paper,margin=24mm]{geometry}
\usepackage{amsmath,amssymb,amsthm,mathtools,booktabs,hyperref}
\newtheorem{theorem}{Theorem}
\newtheorem{lemma}{Lemma}
\newcommand{\Sym}{\operatorname{Sym}}
\newcommand{\ran}{\operatorname{ran}}
\newcommand{\fall}[2]{(#1)_{#2}}
\title{The Bosonic Laughlin Spectral Gap: A Candidate Proof\\[0.5em]\large\textit{An AI-generated research note}}
\author{\normalsize Prepared and curated by Xin Shen\thanks{\textbf{AI-generated candidate proof.}
OpenAI Codex AI agents developed the core argument, verification code,
threshold optimization, and manuscript under Xin Shen's direction and
curation on 7 October 2026. The session identifies the model family as
GPT-6; the exact deployed model identifier was not recorded.
Exact planar and scalar checks and 156 interval block certificates were
replayed successfully. The analytical reduction and its implementation
remain open to independent review; no independent human proof review or
proof-assistant formalization is recorded for v0.1.
See the repository's AI contribution and verification statement.}}
\date{Version 0.1, 7 October 2026; $Q_0=3088$}
\begin{document}
\maketitle
\begin{abstract}
We present a candidate computer-assisted proof of $H_Q^2\ge H_Q/3$ for the full spherical
bosonic $V_0$ interaction at every integer flux $Q\ge3088$, independently
of particle number. Each unordered pair projector has coefficient one.
At flux $Q=2(N-1)$ this gives a gap of at least $1/3$ above the unique
bosonic Laughlin state. The argument combines 138 fixed operator rows,
coherent comparisons resolved by total spin, and 156 parameter-interval
block certificates. Six overlapping flux ranges cover $[3088,\infty)$.
The certificate checks use exact rational arithmetic or outward-rounded
integer intervals.
\end{abstract}

\section*{Introduction}
The bosonic Laughlin state at filling $\nu=1/2$ is an exact zero mode of
the $V_0$ Haldane pseudopotential in the lowest Landau level. The
wavefunction and spherical parent-Hamiltonian construction belong to the
Laughlin--Haldane framework~\cite{Laughlin1983,Haldane1983}. The
second-quantized pseudopotential algebra and its zero modes have also been
studied directly~\cite{OrtizEtAl2013,ChenSeidel2015}. Here we consider the
complete projected interaction on the round sphere, with coefficient one
for each unordered pair projector onto relative angular momentum zero.

The spectral question is whether the positive spectrum stays uniformly
separated from zero as the number of particles and the flux grow.
Rougerie's survey states the Laughlin spectral-gap conjecture and discusses
its geometric variants~\cite[Appendix~A]{Rougerie2019}. For a truncated
bosonic pseudopotential arising from the thin-cylinder regime, Warzel and
Young proved gap bounds uniform in volume and particle
number~\cite{WarzelYoung2023}. The present candidate proof concerns the
full spherical $V_0$ interaction. Writing $H_Q$ for this Hamiltonian at integer flux $Q$, its proposed
estimate is $H_Q^2\ge H_Q/3$ for $Q\ge3088$, uniformly in the particle
number $N$; at $Q=2(N-1)$ this gives the stated gap above the Laughlin zero mode.

The argument adapts the comparison on the full Fock space in
OpenAI's preprint on the fermionic $V_1$ interaction~\cite{OpenAI2026}. Its structure
is a normal-ordered Hamiltonian square, a rotation average of positive
operator rows, finite spin-block comparisons, and transfer from planar
coefficients to finite flux. For the bosonic model, coherent occupation
bounds control the comparison errors before the many-body lift. The
resulting finite inequalities are supplied as frozen data and checked with
exact rational arithmetic and outward-rounded intervals.

\section{Model, normalization, and three-body input}
On the sphere, $Q$ is the integer number of flux quanta. The one-particle LLL
is $U_Q$, the spin-$Q/2$ representation, with normalized orbitals $e_i$,
$0\le i\le Q$. The physical $N$-particle space is $\Sym^N U_Q$.
The pair space $V_Q\subset\Sym^2U_Q$ is the maximal-spin-$Q$ subspace,
the subspace of relative angular momentum zero. Its normalized orbital-sum basis
is $v_p$, $0\le p\le2Q$. With bosonic CCR
$[a_i,a_j^\dagger]=\delta_{ij}$, set
\begin{equation}
 t_p^Q(i,j)=\delta_{i+j,p}
 \sqrt{\frac{\binom Qi\binom Qj}{\binom{2Q}p}},\qquad
 B_p=\frac1{\sqrt2}\sum_{i,j=0}^Q t_p^Q(i,j)a_i a_j,
 \qquad H_Q=\sum_p B_p^\dagger B_p.
\end{equation}
Thus $B_p^\dagger|0\rangle=v_p$ and each unordered pair projector has
coefficient one. The interaction is not orbitally truncated. At
$Q=2(N-1)$ its zero mode is the bosonic Laughlin polynomial
$\prod_{i<j}(u_iv_j-u_jv_i)^2$, where $(u_i,v_i)$ are homogeneous
spinor coordinates of particle $i$ on the sphere; an LLL polynomial is
homogeneous of degree $Q$ in each particle's spinor.

On the plane, $U$ is the Bargmann LLL with normalized orbitals $e_i$,
$i\ge0$. Replace the tensor by
\begin{equation}
 t_p^\infty(i,j)=\delta_{i+j,p}
 \sqrt{\frac{p!}{2^p i!j!}},\qquad H=\sum_{p\ge0}B_p^\dagger B_p.
\end{equation}
All forms below can first be read on finite-particle algebraic Fock vectors.
In each fixed $N$ sector the projected pair Hamiltonian is bounded. The
particle-number-independent inequalities then hold on the appropriate
direct-sum form domain.

For a normalized $k$-particle vector $w$, let $A(w)^\dagger$ be the unique
homogeneous degree-$k$ polynomial in creation operators satisfying
$A(w)^\dagger|0\rangle=w$, and let $A(w)$ be its physical adjoint.
The positive, normal-ordered $k$-body lift is the linear map
\begin{equation}
 \mathcal L_k(|w\rangle\langle w'|)=A(w)^\dagger A(w').
\end{equation}
It takes positive operators on the physical $k$-particle space to positive
forms on all bosonic particle sectors. In an occupation basis $\nu$,
$A(\nu)=\prod_i a_i^{\nu_i}/\sqrt{\nu_i!}$; no extra factorial is present.

Normal ordering gives the exact identity
\begin{equation}\label{eq:normal}
 H_Q^2=H_Q+S_{3,Q}+S_{4,Q},\qquad
 S_{4,Q}=\sum_{p,q}(B_qB_p)^\dagger(B_qB_p)\ge0.
\end{equation}

The auxiliary space $V_Q\otimes U_Q$ contains a pair and a distinguished
third particle. The linear creation map
\begin{equation}
 W_3:V_Q\otimes U_Q\longrightarrow\Sym^3U_Q,
 \qquad v_p\otimes u\longmapsto B_p^\dagger u
\end{equation}
is not an isometry. Couple pair spin $Q$ and single-particle spin $Q/2$.
Let $\Pi_z$ be the auxiliary projection onto spin $J_z=3Q/2-z$,
$0\le z\le Q$, with dimension $d_z=3Q-2z+1$.
The recoupling identities are
\begin{equation}\label{eq:recouple}
 W_3^\dagger W_3=\sum_{z=0}^Q\lambda_{z,Q}\Pi_z,
 \qquad \lambda_{z,Q}=1+2(-1)^z f_{z,Q},\qquad
 f_{z,Q}=\frac{\fall Qz}{\fall{2Q}z},
\end{equation}
and
\begin{equation}
 S_{3,Q}=\mathcal L_3\!\left[
 W_3\Big(2\sum_z(-1)^z f_{z,Q}\Pi_z\Big)W_3^\dagger\right].
\end{equation}
Here $\fall xz=x(x-1)\cdots(x-z+1)$, with $\fall x0=1$.
Let $J:V_Q\otimes U_Q\to U_Q^{\otimes3}$ be the inclusion
putting the pair in the first two tensor slots, and let $P_{\rm sym}$ be the
orthogonal three-particle symmetrizer. Then
$W_3=\sqrt3P_{\rm sym}J$ and
$W_3^\dagger W_3=I+2T$, where $T=J^\dagger(23)J$ is the compressed
transposition of slots two and three. The pair--single spin decomposition is
multiplicity free. A highest vector at deficit $z$ has the form
$\sum_{p=0}^zs_pv_p\otimes e_{z-p}$, with
\begin{equation}
 \frac{s_{p+1}}{s_p}
 =-\sqrt{\frac{(z-p)(Q-z+p+1)}{(p+1)(2Q-p)}},\qquad
 \frac{s_z}{s_0}=(-1)^z\sqrt{f_{z,Q}}.
\end{equation}
Taking the $v_0\otimes e_z$ coefficient after transposition gives
$s_zt_z^Q(0,z)=s_z\sqrt{f_{z,Q}}$; hence $T$ has eigenvalue
$(-1)^zf_{z,Q}$. This proves \eqref{eq:recouple}. Also
$W_3W_3^\dagger=H_{3,Q}$, so its three-body square minus itself is the
displayed three-body coefficient. Direct CCR normal ordering gives
$[B_p,B_q^\dagger]=\delta_{pq}+
2\sum_{i,j,k}t_p^Q(i,k)t_q^Q(j,k)a_j^\dagger a_i$,
and consequently \eqref{eq:normal} on every particle sector by the positive
three-body lift. In particular, $S_3$ is not positive; the $z=1$ auxiliary
block is annihilated by $W_3$.

Fix an odd integer $m\ge3$. The positive magnitude of the negative tail is
\begin{equation}\label{eq:taildef}
 C_{m,Q}=2\sum_{\substack{z\ge m\\z\text{ odd}}}f_{z,Q}\Pi_z,
 \qquad T_{m,Q}=\mathcal L_3(W_3 C_{m,Q}W_3^\dagger).
\end{equation}
In the plane, $f_{z,\infty}=2^{-z}$. The index $z$ there is the degree of
the internal oscillator between pair center of mass and the third particle.

\section{Plane estimate}
\begin{theorem}\label{thm:plane}
For odd $m\ge3$ and every particle number,
\begin{equation}\label{eq:planebound}
 T_m\le c_m(S_4+H),\qquad
 c_m=3(m+1)\left(\frac{m+1}{2m}\right)^m.
\end{equation}
More generally, for $1/2<r<1$ and $\eta>0$,
\begin{equation}
 T_m\le \frac3{1-r}\left(\frac1{2r}\right)^m
 \big(\eta S_4+C_\eta H\big),\qquad
 C_\eta=\max_{n\in\mathbb Z_{\ge0}}
 \left[n-\frac\eta2n(n-1)\right].
\end{equation}
In particular $C_1=1$.
\end{theorem}
\begin{proof}
Use the pair center-of-mass oscillator $u$ and third-particle oscillator $v$. The unitary
oscillator change is
\begin{equation}
 \zeta=\frac{\sqrt2u+v}{\sqrt3},\qquad
 \xi=\frac{-u+\sqrt2v}{\sqrt3}.
\end{equation}
Let $n_R$ count quanta of $\xi$. The auxiliary tail satisfies
\begin{equation}
 C_m=2\,2^{-n_R}P_{\text{odd},\ge m}
 \le2\left(\frac1{2r}\right)^m(I_{\rm CM}\otimes r^{n_R}).
\end{equation}
For normalized oscillator coherent vectors the exact positive representation
is
\begin{equation}
 r^{n_R}=\int\frac{d^2\xi}{\pi r}
 e^{-(1-r)|\xi|^2/r}|\xi\rangle\langle\xi|.
\end{equation}
Its $j$-th diagonal matrix element is $r^j$, obtained by the elementary
Gaussian moment integral. Resolve $I_{\rm CM}$ in coherent states as well.
Write $u=\sqrt2\alpha$, $v=\beta$, where $\alpha,\beta$ are physical
one-particle coherent coordinates. Then
\begin{equation}\label{eq:coherent}
 I_{\rm CM}\otimes r^{n_R}
 =\int\frac{2d^2\alpha\,d^2\beta}{\pi^2r}
 e^{-\kappa|\beta-\alpha|^2}
 |\sqrt2\alpha\rangle_V\langle\sqrt2\alpha|
 \otimes|\beta\rangle\langle\beta|,
 \qquad \kappa=\frac{2(1-r)}{3r}.
\end{equation}
The factor $2$ is the Jacobian of $u=\sqrt2\alpha$.

Let $a_\alpha$ annihilate in the normalized physical coherent orbital,
$B_\alpha=a_\alpha^2/\sqrt2$, and $n_\beta=a_\beta^\dagger a_\beta$.
The pair coherent vector is exactly $B_\alpha^\dagger|0\rangle$.
Positive lifting the integrand in \eqref{eq:coherent} after $W_3$ produces
$B_\alpha^\dagger n_\beta B_\alpha$. Since the coherent pair belongs to $V$,
\begin{equation}\label{eq:occupation}
 H\ge B_\beta^\dagger B_\beta=\tfrac12n_\beta(n_\beta-1),\qquad
 n_\beta\le\eta H+C_\eta I.
\end{equation}
The latter inequality acts on the output sector after
$B_\alpha$ removes two particles. The Gaussian $\beta$ marginal is
$3/[2(1-r)]$. The pair coherent resolution gives the exact identities
\begin{equation}\label{eq:frames}
 \int\frac{2d^2\alpha}{\pi}B_\alpha^\dagger B_\alpha=H,\qquad
 \int\frac{2d^2\alpha}{\pi}B_\alpha^\dagger H B_\alpha
 =\sum_p B_p^\dagger H B_p=S_4.
\end{equation}
Pair annihilators commute, which explains the last equality. Consequently
\begin{equation}
 \mathcal L_3[W_3(I_{\rm CM}\otimes r^{n_R})W_3^\dagger]
 \le\frac3{2(1-r)}(\eta S_4+C_\eta H).
\end{equation}
Combine with the auxiliary spectral bound. At $\eta=1$, minimize over $r$;
the minimum occurs at $r=m/(m+1)$ and gives \eqref{eq:planebound}.
\end{proof}

\section{Spherical estimate with the same asymptotic constant}
\begin{theorem}\label{thm:sphere}
Fix odd $m\ge3$ and $\varepsilon>0$. There exists a flux threshold
$Q_{m,\varepsilon}$ independent of particle number such that
\begin{equation}\label{eq:spherebound}
 T_{m,Q}\le(c_m+\varepsilon)(S_{4,Q}+H_Q)
 \quad(Q\ge Q_{m,\varepsilon})
\end{equation}
on all bosonic particle sectors, with the same $c_m$ as in
Theorem~\ref{thm:plane}. The threshold is existential in this argument.
\end{theorem}
\begin{proof}
Let $d=Q+1$, $d_a=2Q+1$. Sphere measures have total mass one. Fix
$0<a<1/3$, $M=\lfloor aQ\rfloor$, and introduce the positive auxiliary kernel
\begin{equation}
 A_{M,Q}=d_ad\int d\Omega_\alpha d\Omega_\beta\,
 s(\alpha,\beta)^M
 |\alpha_{\rm pair}\rangle\langle\alpha_{\rm pair}|
 \otimes|\beta\rangle\langle\beta|,
 \quad s=\cos^2(\theta_{\alpha\beta}/2).
\end{equation}
These are normalized spin-coherent vectors. Rotation invariance gives
$A_{M,Q}=\sum_z\tau_{z,Q}\Pi_z$.

Fix $\alpha$ at the north pole. A beta integral gives the diagonal coefficient
\begin{equation}
 b_{j,Q}=d\binom Qj\int_0^1s^{M+Q-j}(1-s)^j ds
 =\frac{Q+1}{M+Q+1}\frac{\fall Qj}{\fall{M+Q}j}.
\end{equation}
Tracing in the coupled spin block yields
\begin{equation}\label{eq:fulltau}
 \tau_{z,Q}=\frac{d_a}{d_z}\sum_{j=z}^Q b_{j,Q}\gamma_{z,j,Q},\quad
 \gamma_{z,j,Q}=
 \frac{\fall{2Q}z}{\fall{3Q-z+1}z}
 \binom jz\frac{\fall{Q-z}{j-z}}{\fall{3Q-2z}{j-z}}.
\end{equation}
Here $\gamma$ is the squared Clebsch--Gordan coefficient for a highest pair
and a single orbital lowered $j$ times. To derive it, normalize the highest
coupled-spin vector using
\begin{equation}
 \frac{|c_p|^2}{|c_{p-1}|^2}
 =\frac{(z-p+1)(Q-z+p)}{p(2Q-p+1)}.
\end{equation}
The terminating Chu--Vandermonde identity gives
$|c_0|^2=\fall{2Q}z/\fall{3Q-z+1}z$. After total lowering $j-z$ times,
only lowering the single-particle component can retain pair highest weight;
the normalization ratio gives the remaining binomial and falling-factorial
factors in \eqref{eq:fulltau}.

We first need an all-$z$ coarse bound. Retaining only the positive $j=z$
term in \eqref{eq:fulltau}, with $\rho=3(1+a)/4<1$, gives
\begin{equation}\label{eq:coarse}
 \frac{2f_{z,Q}}{\tau_{z,Q}}\le3(1+a)\rho^z.
\end{equation}
Indeed the ratio before the factor $2$ is at most
\begin{equation}
 \frac{d_z}{d_a}\frac{M+Q+1}{d}
 \prod_{i=0}^{z-1}
 \frac{(M+Q-i)(3Q-z+1-i)}{(2Q-i)^2}.
\end{equation}
The prefactor is at most $\tfrac32(1+a)$. Since $i\le z-1$ and $M\le aQ$,
each product factor is at most
$(1+a-x)(3-2x)/(2-x)^2\le\rho$, $x=i/Q\in[0,1)$.
The last inequality is equivalent to
$x[(8-4a)-(5-3a)x]\ge0$.

For each fixed $z$, the full sum in \eqref{eq:fulltau} has the limit
\begin{equation}\label{eq:taulimit}
 \tau_{z,Q}\longrightarrow r^{z+1},\qquad r=\frac2{2+3a}.
\end{equation}
For fixed $z,j$, the limiting summand uses
$b_{j,Q}\to(1+a)^{-j-1}$ and
$\gamma_{z,j,Q}\to\binom jz(2/3)^z(1/3)^{j-z}$,
while $d_a/d_z\to2/3$. The binomial-series identity sums these limits.
To justify exchanging limit and sum, extend summands by zero for $j>Q$.
For $Q\ge2/a$, $M\ge aQ/2$,
\begin{equation}
 0\le b_{j,Q}\le(1+a/2)^{-j},\quad
 0\le\gamma_{z,j,Q}\le1,\quad d_a/d_z<2.
\end{equation}
This provides a summable majorant independent of $Q$.

Now fix $m$ and choose $a=2/(3m)$, so $r=m/(m+1)>1/2$. The fixed-$z$ ratio
limit is $(2/r)(1/(2r))^z$, decreasing on odd $z\ge m$. Pick a finite
$L$ so large that the right side of \eqref{eq:coarse} for $z>L$ is at most
its limiting maximum at $z=m$. On the finite set $m\le z\le L$ use
\eqref{eq:taulimit}. For any $\delta>0$, sufficiently large $Q$ therefore
gives the whole-space inequality
\begin{equation}
 C_{m,Q}\le\left[\frac2r\left(\frac1{2r}\right)^m+\delta\right]A_{M,Q}.
\end{equation}

Finally lift the coherent kernel. Equation~\eqref{eq:occupation} holds on
the sphere as well. The exact beta marginal is $d/(M+1)$, while the pair
coherent resolution has factor $d_a$. Thus
\begin{equation}
 \mathcal L_3(W_3A_{M,Q}W_3^\dagger)
 \le\frac{d}{M+1}(S_{4,Q}+H_Q).
\end{equation}
Since $d/(M+1)\to1/a$, the resulting scalar tends to
$(2/(ra))(1/(2r))^m=c_m$. Choose $\delta$ small and then $Q$ sufficiently
large. All comparisons precede lifting and their constants are independent
of $N$, proving \eqref{eq:spherebound}.
\end{proof}


\section{Absorbing finite three- and four-body comparison errors}
Write $H_{r,Q}$ for $H_Q$ restricted to the physical $r$-particle space,
where $r=3,4$. Let $\Pi^{(r)}_{\le D,Q}$ be its total-spin projection onto
$J=rQ/2-z$, $0\le z\le D$, including all multiplicities. Define the
many-body positive forms
\begin{equation}
 \mathcal R_{r,D,Q}:=\mathcal L_r(H_{r,Q}\Pi^{(r)}_{\le D,Q}).
\end{equation}
The product is positive because the Hamiltonian is rotationally invariant.
The cutoff is in the few-body coefficient, not in the physical many-body
Hamiltonian or in particle occupation.

\begin{lemma}\label{lem:finite-errors}
For all sufficiently large $Q$, uniformly in total particle number,
\begin{equation}\label{eq:error-bounds}
 \mathcal R_{3,8,Q}\le35(S_{4,Q}+H_Q),\qquad
 \mathcal R_{4,16,Q}\le2100(2S_{4,Q}+H_Q).
\end{equation}
Both estimates hold, in particular, for every integer $Q\ge10^{50}$,
as shown in Appendix~\ref{app:explicit-coherent-threshold}. Their thresholds
depend only on the displayed fixed cutoffs.
\end{lemma}
\begin{proof}
The three-body claim follows from the positive auxiliary operator
$A_{M,Q}$ of the preceding section. On each retained spin block its
eigenvalue tends to $r^{z+1}$, $r=2/(2+3a)$. Since only $z\le D$ is now
retained, positivity outside this finite set and finite-block convergence
give
\begin{equation}
 \Pi^{\rm aux}_{\le D,Q}\le(r^{-(D+1)}+o(1))A_{M,Q}.
\end{equation}
After positive lifting, the limiting constant is
$a^{-1}(1+3a/2)^{D+1}$. Taking $D=8,a=1/12$ gives
$12(9/8)^9=34.63809\ldots<35$.

For the four-body claim, use the wide pair map
\begin{equation}
 W_4:V_Q\otimes\Sym^2U_Q\longrightarrow\Sym^4U_Q,
 \qquad W_4(v_p\otimes u)=B_p^\dagger u.
\end{equation}
It satisfies $W_4W_4^\dagger=H_{4,Q}$. Its domain differs from the narrower
pair--pair space that represents $S_4$. Equivariance gives
$W_4\Pi^{\rm aux}_{\le D,Q}W_4^\dagger
=H_{4,Q}\Pi^{(4)}_{\le D,Q}$.

First establish an occupation bound for possibly nonorthogonal coherent
orbitals $\beta,\gamma$. Put $A=a_\gamma a_\beta$ and
$s=\langle\beta,\gamma\rangle$, with $|s|\le1$. Normal ordering and
$(n_\beta-n_\gamma)^2\ge0$ imply
\begin{align}
 A^\dagger A
 &=\tfrac12(n_\beta n_\gamma+n_\gamma n_\beta)
       -\operatorname{Re}(s a_\beta^\dagger a_\gamma)\notag\\
 &\le\tfrac12(n_\beta^2+n_\gamma^2)
       +\tfrac12(n_\beta+n_\gamma)
 \le4H_Q+2I.\label{eq:two-spectators}
\end{align}
For the middle inequality, quadratic-form Cauchy--Schwarz bounds the real
cross term by $|s|(n_\beta+n_\gamma)/2$. The last step uses
$n_x^2\le2H_Q+n_x$ and $n_x\le H_Q+I$ for each coherent $x$.

Set $M=\lfloor aQ\rfloor$ and introduce the positive invariant operator
on $V_Q\otimes\Sym^2U_Q$
\begin{align}
 A^{(4)}_{M,Q}=d_ad^2\int &d\Omega_\alpha d\Omega_\beta d\Omega_\gamma\,
 [s(\alpha,\beta)s(\alpha,\gamma)]^M
 |\alpha_{\rm pair}\rangle\langle\alpha_{\rm pair}|\notag\\[-2pt]
 &\otimes P_{\rm sym}^{(2)}|\beta\gamma\rangle
       \langle\beta\gamma|P_{\rm sym}^{(2)}.
\end{align}
The isometry from symmetric tensors to bosonic Fock space sends
$P_{\rm sym}^{(2)}|\beta\gamma\rangle$ to
$a_\beta^\dagger a_\gamma^\dagger|0\rangle/\sqrt2$.
Consequently the lifted coherent integrand is
$\tfrac12B_\alpha^\dagger A^\dagger A B_\alpha$, and
\eqref{eq:two-spectators}, the exact two spectator marginals, and the pair
frame give
\begin{equation}\label{eq:four-kernel-lift}
 \mathcal L_4(W_4A^{(4)}_{M,Q}W_4^\dagger)
 \le\left(\frac{d}{M+1}\right)^2(2S_{4,Q}+H_Q).
\end{equation}

We next compute the limit of every fixed auxiliary spin block, including
its multiplicities. Choose continuous normalized highest-weight bases at
deficit $z$, and normalized total-lowering descendants numbered by $n$.
Such bases exist for large $Q$: the finite-weight raising matrices converge
to the oscillator lowering matrix, which has full row rank. More explicitly,
the auxiliary weight-$z$ basis consists of $v_p\otimes|j,k\rangle_{\rm sym}$,
$p+j+k=z$, $j\le k$. For $Q\ge z$ its dimension is fixed, and the
highest-weight multiplicity is $\lfloor z/2\rfloor+1$. After division by
$\sqrt Q$, the raising matrix tends to $2a_\zeta$, where $a_\zeta$ annihilates
the total center-of-mass oscillator introduced below. It is surjective from
weight $z$ to weight $z-1$. For a full-row-rank matrix $T$, the kernel
projector is $I-T^\dagger(TT^\dagger)^{-1}T$; this formula proves convergence
of these kernel projections and permits convergent orthonormal bases.
The case $z=0$ is the one-dimensional highest line and needs no inverse.
Normalized spherical descendants then converge, for each fixed descendant
number $n$, to $(a_\zeta^\dagger)^n/\sqrt{n!}$ applied to the limiting
highest vectors.
Fix the pair at the north pole. The two-spectator operator is the restriction of
$R_Q\otimes R_Q$ to symmetric tensors, where $R_Q$ is diagonal with
eigenvalues $b_{j,Q}$ from the preceding section. Schur averaging expresses
each multiplicity-matrix entry as
\begin{equation}
 \frac{d_a}{4Q-2z+1}\sum_{n\ge0}
 \sum_{\substack{j\le k\\j+k=z+n}}
 b_{j,Q}b_{k,Q}\,\overline{c_{b,jk,n,Q}}c_{c,jk,n,Q}.
\end{equation}
Here $c_{b,jk,n,Q}$ is the pair-highest component of the normalized
descendant of highest vector $b$; the spectator basis is orthonormal, also
when $j=k$. The squared sum of these components is at most one. For
$Q\ge2/a$, the inner sum is therefore bounded in absolute value by
$q^{z+n}$, where $q=(1+a/2)^{-1}<1$. The dimension ratio is bounded.
This summable majorant justifies dominated convergence of the entire
descendant sum.

To identify that limit, use the unitary Jacobi-oscillator change
\begin{equation}
 u=\sqrt2\alpha,\quad v=\beta,\quad w=\gamma,\qquad
 \zeta=\frac{\sqrt2u+v+w}{2},\quad
 \xi_1=\frac{v-w}{\sqrt2},\quad
 \xi_2=\frac{u}{\sqrt2}-\frac{v+w}{2}.
\end{equation}
In this limit the pair-north seed is the positive trace-class operator
\begin{equation}
 \sigma=|0\rangle_u\langle0|\otimes
 \left[R\otimes R\right]_{\rm sym},\qquad
 R=(1+a)^{-(n+1)}
   =\int\frac{d^2\beta}{\pi}e^{-a|\beta|^2}|\beta\rangle\langle\beta|,
\end{equation}
where $n$ is the one-oscillator occupation operator in either spectator
factor. On relative degree $z$, the full descendant sum is the restriction
of the center-of-mass partial trace of $\sigma$, with the prefactor
$d_a/(4Q-2z+1)\to1/2$.
Let $D_\zeta(h)=\exp(h a_\zeta^\dagger-\bar h a_\zeta)$ be
the Weyl displacement on the center-of-mass factor. Simultaneously
displacing $(u,v,w)$ by $(\sqrt2\alpha,\alpha,\alpha)$ displaces only
$\zeta$, by $2\alpha$. The Weyl twirling identity for the positive
trace-class seed therefore gives, as a weak operator integral,
\begin{equation}\label{eq:twirl-audit}
 \int\frac{2d^2\alpha}{\pi}
 D_\zeta(2\alpha)\sigma D_\zeta(2\alpha)^\dagger
 =\frac12 I_{\rm CM}\otimes\operatorname{Tr}_{\rm CM}\sigma.
\end{equation}
The factor $1/2$ follows from $d^2(2\alpha)=4d^2\alpha$ and agrees with
the Schur dimension ratio. Inserting the coherent representation of $R$
into the left side of \eqref{eq:twirl-audit} gives the measure
$2d^2\alpha\,d^2\beta\,d^2\gamma/\pi^3$ with weight
$e^{-a(|\beta-\alpha|^2+|\gamma-\alpha|^2)}$.
The twirling identity can also be checked in center-of-mass position
coordinates $x,y$. For $h=(q+ip)/\sqrt2$, conjugation by $D_\zeta(h)$
sends the operator-valued kernel to
$e^{ip(x-y)}\sigma(x-q,y-q)$. Integrating $dq\,dp/(2\pi)$ first gives
$\delta(x-y)$ and then the partial trace. Here the Gaussian, positive
trace-class seed permits this identity in the weak sense.
The substitution $u=\sqrt2\alpha$ cancels this factor two. Thus the
limiting coherent kernel has weight
$\exp[-a|\xi_1|^2-2a|\xi_2|^2]$, and hence is exactly
\begin{equation}
 A^{(4)}_{\infty}=I_{\rm CM}\otimes r_1^{n_1+1}r_2^{n_2+1},\qquad
 r_1=(1+a)^{-1},\quad r_2=(1+2a)^{-1}.
\end{equation}
Here $n_i=a_{\xi_i}^\dagger a_{\xi_i}$ counts quanta of relative oscillator
$\xi_i$, for $i=1,2$.
Spectator symmetry restricts $n_1$ to even values. At fixed deficit
$n_1+n_2=z$, the least eigenvalue is $r_1r_2^{z+1}$. Thus on the finite
set $z\le D$, continuity gives
\begin{equation}
 \Pi^{\rm aux}_{\le D,Q}
 \le\big[(r_1r_2^{D+1})^{-1}+o(1)\big]A^{(4)}_{M,Q}.
\end{equation}
The inequality holds on the whole auxiliary space because invariance
prohibits cross terms between distinct spins and the right side is positive
outside the retained blocks.
Combining with \eqref{eq:four-kernel-lift}, the limiting scalar is
\begin{equation}
 C_{4,D}(a)=\frac{(1+a)(1+2a)^{D+1}}{a^2}.
\end{equation}
At $D=16,a=1/16$, this equals $272(9/8)^{17}=2014.47455\ldots<2100$.
All constants and thresholds were obtained before the many-body lift and
are independent of total particle number.
\end{proof}

\section{Frozen finite certificate and its exact verification}
We use rows with net orbital weight $t\in\{0,1,2,3\}$,
\begin{equation}\label{eq:rows}
 F_\rho=\lambda_\rho B_{t_\rho}
 +\sum_{\substack{p,j\ge0\\t_\rho\le p+j\le8}}
 \alpha^\rho_{pj}a_{p+j-t_\rho}^\dagger a_jB_p.
\end{equation}
All coefficients are fixed real algebraic numbers, independent of $Q,N$.
They are defined exactly by the companion file
\path{exact_shifted_certificate_rows.json}, as specified below. The four
channels contain $31,37,36,34$ rows, respectively, totaling $138$.
Set
\begin{equation}
 K_Q=(2Q+1)\int\sum_\rho F_\rho(g)^\dagger F_\rho(g)\,dg
 =\eta H_Q+K_{3,Q}+K_{4,Q}\ge0,\qquad
 \eta=\sum_\rho\lambda_\rho^2.
\end{equation}
Haar measure has mass one. Normal ordering terminates at body order four.
Here $g\in SU(2)$ and $F_\rho(g)=U(g)F_\rho U(g)^\dagger$, with $U(g)$
the rotation representation on the direct sum of all particle sectors.
Before averaging, the three-body coefficient has orbital-weight support at
most $8$, and the four-body coefficient at most $16$. After $SU(2)$
averaging their spin-deficit supports are $z\le8$ and $z\le16$, because
an irrep of deficit $z$ contains no orbital-weight space $M<z$.
The averaged comparison coefficients extend over all descendants of the retained
irreps; $H_Q$ remains untruncated.
Suppress the row index and put $i=p+j-t$, $i'=p'+j'-t$. For a real row, its unaveraged
normal-ordered terms beyond $\lambda^2B_t^\dagger B_t$ are
\begin{align}
 K_3^{\rm row}
 &=\lambda\sum_{p,j}\alpha_{pj}
    (B_t^\dagger a_i^\dagger a_jB_p+\mathrm{h.c.})\notag\\
 &\quad+\sum_{p,j,p',j'}\alpha_{pj}\alpha_{p'j'}\delta_{ii'}
     B_p^\dagger a_j^\dagger a_{j'}B_{p'},\notag\\
 K_4^{\rm row}
 &=\sum_{p,j,p',j'}\alpha_{pj}\alpha_{p'j'}
     B_p^\dagger a_j^\dagger a_{i'}^\dagger a_i a_{j'}B_{p'}.
 \label{eq:normal-rows-audit}
\end{align}
The three-body weights are $p+j\le8$ and the four-body weights are
$(p+j)+(p'+j')-t\le16$. Every term retains a pair annihilator on the right
and a pair creator on the left. These are exact CCR identities on each sphere.

Let $\mathsf K_{r,Q}$ denote the physical $r$-particle coefficient whose
lift is $K_{r,Q}$. Let
$\mathsf S_{3,Q}=H_{3,Q}^2-H_{3,Q}$ and let $\mathsf S_{4,Q}$ be the
physical four-body coefficient of $S_{4,Q}$. Define
\begin{align}
 E_{3,Q}&=\Pi^{(3)}_{\le8,Q}\mathsf S_{3,Q}\Pi^{(3)}_{\le8,Q}
                    -\mathsf K_{3,Q},\notag\\
 E_{4,Q}&=\Pi^{(4)}_{\le16,Q}
  \big[(1-c_9-\mu)\mathsf S_{4,Q}-\mathsf K_{4,Q}\big]
                  \Pi^{(4)}_{\le16,Q},\label{eq:finiteE}\\
 c_9&=30(5/9)^9=\frac{19531250}{129140163},\qquad \mu=\frac1{50}.
\end{align}
Outside the four-body cutoff the comparison coefficient is a
positive multiple of $\mathsf S_{4,Q}$.

The finite certificate is evaluated in planar limiting relative blocks.
If $L=\sum_{j\ge1}\sqrt j\,a_{j-1}^\dagger a_j$ on the $r$-particle
space, then $[L,L^\dagger]=r$. A highest (relative) vector at degree $z$
has $L\psi=0$, and its normalized degree-$M$ descendant is
$(L^\dagger)^{M-z}\psi/\sqrt{r^{M-z}(M-z)!}$.
Denote this isometric map from the relative highest-weight space to the
degree-$M$ physical space by $U_{r,z}^{(M)}$. In the next formula,
$\mathsf K_{r,M}$ is the unaveraged normal-ordered body coefficient
restricted to total orbital degree $M$; the overbar denotes its averaged
relative block.
Schur averaging on the sphere, followed by the fixed-block limit, gives
\begin{equation}\label{eq:average-formula}
 \overline{\mathsf K}_{r,z}
 =\frac2r\sum_{M\ge z}
 (U_{r,z}^{(M)})^\dagger\mathsf K_{r,M}U_{r,z}^{(M)}.
\end{equation}
The sum is finite. The factor is
$(2Q+1)/(rQ-2z+1)\to2/r$. This also follows from the oscillator
center-of-mass partial trace. Two-body averaging is exactly $\eta H$.

To make the verification entirely rational, introduce canonical polynomial
operators
\begin{equation}
 x_j=\sqrt{j!}\,a_j^\dagger,\qquad \partial_j=a_j/\sqrt{j!}.
\end{equation}
They satisfy $[\partial_i,x_j]=\delta_{ij}$, but are not physical adjoints:
$\partial_j^*=x_j/j!$ and $x_j^*=j!\partial_j$.
For the unnormalized physical occupation basis
$|A]_x=\prod_jx_j^{n_j}|0\rangle$, the metric is the positive integer
$g_A=\prod_jn_j!(j!)^{n_j}$. On the plane,
\begin{equation}
 B_p=\kappa_p\widehat B_p,\qquad
 \kappa_p^2=\frac{p!}{2^{p+1}},\qquad
 \widehat B_p=\sum_{i+j=p}\partial_i\partial_j.
\end{equation}
The frozen JSON integers divided by $10^{12}$ define the scaled row
coefficients
\begin{equation}\label{eq:scaledrows}
 y_0=\lambda\kappa_t,\qquad
 y_{pj}=\alpha_{pj}\kappa_p\sqrt{j!/i!},\qquad i=p+j-t.
\end{equation}
Equations \eqref{eq:rows} and \eqref{eq:scaledrows} uniquely define the
same fixed coefficients used on every sphere.

For the finite verification,
compute $F^*F$ in the $2,3,4$-particle sectors with the metric just given.
Here $F^*F$ is summed over all rows, and $\mathcal L_r^N$ denotes
$\mathcal L_r$ restricted to the physical $N$-particle sector.
Extract the normal-ordered coefficients recursively:
\begin{align}
 \mathsf K_2&=(F^*F)|_{N=2},\notag\\
 \mathsf K_3&=(F^*F)|_{N=3}-\mathcal L_2^3(\mathsf K_2),\notag\\
 \mathsf K_4&=(F^*F)|_{N=4}-\mathcal L_2^4(\mathsf K_2)
                                      -\mathcal L_3^4(\mathsf K_3).
\end{align}
The coefficient at body order $r$ is recovered because its normalized lift
restricts to itself on $r$ particles. There are no higher body orders in
$F^*F$. The integer CM operators are
$L=\sum_{j\ge1}j x_{j-1}\partial_j$ and
$L^\dagger=\sum_{j\ge0}x_{j+1}\partial_j$.
Rational row reduction gives a highest-weight basis. To use
\eqref{eq:average-formula} without square roots, use unnormalized descendants
and divide their quadratic forms by $r^{M-z}(M-z)!$.

With $\delta=10^{-6}$, the exact rational tests give
\begin{equation}\label{eq:exactfinite}
 E_{3,\infty}+\delta H_{3,\infty}\Pi^{(3)}_{\le8,\infty}\ge0,
 \qquad
 E_{4,\infty}+\delta H_{4,\infty}\Pi^{(4)}_{\le16,\infty}\ge0.
\end{equation}
All $9+17=26$ relative blocks were tested, including zero-dimensional ones.
Positive-pivot rational Schur elimination establishes semidefiniteness;
no floating eigenvalue tolerance enters. In every block the rank equals
the rank of $H_r$, and the exact kernel of $H_r$ is preserved. Thus each
shifted matrix is strictly positive on the physical image of $H_r$.
The exact two-body cost is
\begin{equation}\label{eq:eta}
 \eta=\frac{764191703750823432978457}{1500000000000000000000000}
       <\frac{51}{100}.
\end{equation}
The full rational pivots, ranks, and zero-kernel checks are in
\path{exact_shifted_certificate_results.json}.

Only three adjacent files are needed to replay the proof of
\eqref{eq:exactfinite}: the frozen JSON rows,
\path{exact_relative_blocks.py}, and
\path{exact_shifted_certificate.py}. Running
\begin{quote}\ttfamily python -S exact\_shifted\_certificate.py\end{quote}
uses Python integers and
\texttt{fractions.Fraction} to verify the frozen rows. The optional proposal
mode is separate from the default verifier. The verifier returns a nonzero exit status on failure.

\section{Transferring the finite certificate to spheres}
Fix $r=3,4$ and $z\le16$. In the sphere weight-$z$ occupation basis, all
pair coefficients and the rescaled raising/lowering matrices converge to
their planar counterparts. For $Q\ge z$ the finite weight spaces have
fixed dimensions. Dividing the spherical lowering operator by $\sqrt Q$
gives $L^\dagger$ in the limit, and normalizing descendants gives
\eqref{eq:average-formula}. The finite normal-order supports make this a
finite sum for the certificate.

The relevant zero spaces also have exactly stable dimensions. In the
bosonic LLL, $H_{r,Q}\psi=0$ if and only if the polynomial vanishes at
every pair coincidence. Symmetry makes this vanishing at least quadratic,
so the polynomial is divisible by
\begin{equation}
 \Delta_r^2=\prod_{i<j}(u_iv_j-u_jv_i)^2.
\end{equation}
Conversely this divisibility annihilates every $V_0$ pair channel.
The factor is an $SU(2)$ invariant and reduces each single-particle degree
by $2(r-1)$. Let $p_r(z)$ count partitions of $z$ into at most $r$ parts,
with $p_r(z)=0$ for $z<0$, and put $d_r(z)=p_r(z)-p_r(z-1)$.
For $Q\ge z$, the relative physical dimension is $d_r(z)$ and the relative
zero-space dimension is $d_r(z-r(r-1))$. Thus the active rank is exactly
\begin{equation}\label{eq:stable-rank}
 d_r(z)-d_r(z-r(r-1)),
\end{equation}
the same as on the plane.
Indeed multiplication by $\Delta_r^2$ is an injective $SU(2)$ intertwiner
from the residual symmetric polynomial space with single-particle degree
$Q'=Q-2(r-1)$. It shifts total orbital degree by $r(r-1)$, while preserving
total spin. If $z\ge r(r-1)$ and $Q\ge z$, the residual weight
$z'=z-r(r-1)$ obeys $z'\le Q'$, so the orbital cap is inactive in both
weight spaces used to compute the highest-weight multiplicity. If
$z<r(r-1)$ there are no zero modes in this block.

The unprojected normal-ordered coefficients defining $E_{3,Q},E_{4,Q}$
have a pair annihilator on the right and a pair creator on the left.
The spin projections preserve $\ker H_{r,Q}$ because $H_{r,Q}$ is
rotationally invariant. Thus $E_{r,Q}$ annihilates this kernel exactly.
The kernel may rotate with $Q$, but the ranks in
\eqref{eq:stable-rank} are constant. Identify the highest spaces by convergent
orthonormal frames. In one block write $H_Q,E_Q$ for the corresponding
matrices, $P_Q$ for the orthogonal projection onto $\ran H_Q$, and
$G_Q=E_Q+\delta H_Q$ with the planar allowance $\delta=10^{-6}$.
Stable rank and norm convergence imply $P_Q\to P_\infty$: the positive
spectrum of $H_\infty$ is separated from zero, and its spectral cluster
projection is continuous; rank equality rules out additional small positive
eigenvalues. Exact kernel preservation gives $G_Q=P_QG_QP_Q$.
On any nonzero active block let $g>0$ be the least eigenvalue of
$G_\infty$ restricted to $\ran H_\infty$, whose strict positivity was
verified exactly. For a unit vector $x\in\ran H_Q$,
\begin{equation}\label{eq:transfer-margin}
 \langle x,G_Qx\rangle
 \ge g\left(1-\|P_Q-P_\infty\|^2\right)
       -\|G_Q-G_\infty\|.
\end{equation}
This follows from $G_\infty\ge gP_\infty$ and
$\|(I-P_\infty)x\|\le\|P_Q-P_\infty\|$.
For all sufficiently large $Q$, both norms are small enough to make the
right side positive. A zero-rank block is identically zero by kernel
preservation. Thus $G_Q\ge0$ on the entire block, including its rotating
zero space.
In particular the following slightly weaker bounds hold with a fixed allowance
$\delta_Q=2\times10^{-6}$:
\begin{equation}\label{eq:sphere-head}
 E_{3,Q}\ge-\delta_QH_{3,Q}\Pi^{(3)}_{\le8,Q},\qquad
 E_{4,Q}\ge-\delta_QH_{4,Q}\Pi^{(4)}_{\le16,Q}.
\end{equation}
There are only finitely many blocks. The threshold depends on their
matrices and the fixed certificate and is uniform in the total particle number
to which these inequalities are lifted. Appendix~\ref{sec:explicit-head}
gives a quantitative proof using exact polynomial coordinates and an
additional finite planar half-gap check; it establishes these same
inequalities for every integer $Q\ge10^{50}$.

\section{The resulting Fock-space inequality}
\begin{theorem}\label{thm:gap}
For the full spherical bosonic $V_0$ interaction with coefficient one per
unordered pair projector, the explicit integer threshold $Q_0=3088$
works independently of particle number:
\begin{equation}\label{eq:gap}
 H_Q^2\ge\frac13H_Q\qquad(Q\ge Q_0)
\end{equation}
on the entire bosonic LLL Fock space, as a quadratic-form inequality.
For $Q=2(N-1)\ge Q_0$,
\begin{equation}
 H_{N,Q}\ge\frac13(I-P_{\mathrm L,N}),
\end{equation}
where $P_{\mathrm L,N}$ projects onto the normalized bosonic Laughlin state.
The plane interaction also satisfies $H^2\ge H/3$ on every particle sector.
\end{theorem}
\begin{proof}
For $Q\ge3088$, select a row of Table~\ref{tab:optimized-budget}
whose flux interval contains $Q$. Write $L_3,L_4$ for its two exact
rational error budgets and $g$ for its tail improvement. The row ending
at a finite flux $q_+$ uses $g=391/(100q_+)$; the last row uses $g=0$.
Put $\nu=10^{-10}$ and retain $\mu=1/50$.

Appendix~\ref{app:optimized-thresholds} supplies a direct certificate on
the entire parameter interval $0\le1/Q\le1/q_-$ for each row, where $q_-$
is its lower endpoint. Its comparison error in spin block $z$ is the exact
rational allowance $\delta_{r,z}=L_r/\overline C_{r,z}(q_-)$.
The degree-resolved coherent comparison in
Appendix~\ref{app:optimized-coherent} absorbs these errors with cost
$L_3(S_{4,Q}+H_Q)$ and $L_4(2S_{4,Q}+H_Q)$, respectively.
The common positive coherent kernel dominates all its orthogonal spin
blocks, so the cost is the maximum over $z$. Outside the retained blocks, the only negative contribution
is the three-body tail. Thus the positive lifts obey
\begin{align}
 S_{3,Q}-K_{3,Q}&\ge-L_3(S_{4,Q}+H_Q)-T_{9,Q},\notag\\
 (1-c_9-\mu)S_{4,Q}-K_{4,Q}&\ge-L_4(2S_{4,Q}+H_Q).
\end{align}
The exact tail improvement, valid for all $Q\ge3000$, is
\begin{equation}
 T_{9,Q}\le(c_9+\nu-391/(100Q))(S_{4,Q}+H_Q)
          \le(c_9+\nu-g)(S_{4,Q}+H_Q).
\end{equation}
For the last flux interval the second inequality follows by dropping the
negative correction. Substituting into $H_Q^2=H_Q+S_{3,Q}+S_{4,Q}$ and
$K_Q=\eta H_Q+K_{3,Q}+K_{4,Q}$ gives
\begin{align}
 H_Q^2&\ge K_Q+A H_Q+B S_{4,Q},\label{eq:finalbudget}\\
 A&=1-\eta-c_9-\nu+g-L_3-L_4,\notag\\
 B&=\mu-\nu+g-L_3-2L_4.\notag
\end{align}
For every table row, exact rational comparison gives $A>1/3$ and
$B>1/100$. The smallest margin $A-1/3$ is greater than
$42/10^8$. Since $K_Q,S_{4,Q}\ge0$, this proves \eqref{eq:gap}.
Every estimate is on fixed few-body coefficients before the positive
many-body lift, so the flux threshold is independent of total $N$.

At flux $Q=2(N-1)$ the zero-mode factor $\Delta_N^2$ exhausts the
single-particle degree, so the residual symmetric polynomial is constant.
The kernel is therefore the one-dimensional Laughlin line, and the claimed
ground-state gap follows by spectral calculus. In particle sectors with
no zero mode, \eqref{eq:gap} means that the spectrum avoids $(0,1/3)$;
it does not assert the spacing between the two lowest positive levels.

Finally fix a plane particle number and total orbital degree. For all large
$Q$ that complete finite-dimensional sector embeds into the sphere sector
and its Hamiltonian matrices converge to the plane matrices. Passing
\eqref{eq:gap} to this finite-matrix limit gives the same constant.
The direct sum over all degree sectors and then all particle numbers proves
the plane form inequality.
\end{proof}

\section{Reproducibility and scope}
The frozen certificate was replayed in a fresh Python process with site
packages disabled. All 26 blocks passed the exact matrix, rank, and kernel
checks.
An independent implementation reconstructing normalized occupation-space
operators from the frozen rows agreed with the exact two-body cost and all
blocks. Separate checks verified the lift normalization
$\mathcal L_r^N(I)=\binom NrI$, normal ordering through five-particle
sectors, CM descendant metrics, and the finite-$Q$ recoupling formula.
The coherent-frame constants and their normalizations were also
cross-checked.

A separate implementation reconstructed the four-body auxiliary kernel
directly from the total raising matrix. For $a=1/16$, all
descendants contributing to the pair-north Schur sum were included at
$Q=32,64,128$, for $z=0,1,2,4,8,16$. At $z=16$ the smallest eigenvalues
were respectively $0.101477757$, $0.114709929$, and $0.120986698$,
approaching the predicted limit $0.127080285$. Setting the coherent weight
exponent to zero recovered the identity frame with maximum residual
$1.1\times10^{-13}$. Tensor--Fock normalization and the coherent double
integral were also checked independently. These checks used
floating-point arithmetic.

Separately, the frozen rows were evaluated with the actual spherical pair
coefficients and normalized $SU(2)$ descendants. For each retained block
define the energy-relative residual
\begin{equation}
 e_{r,z,Q}:=\lambda_{\min}\!\left[
 H_{r,z,Q}^{-1/2}E_{r,z,Q}H_{r,z,Q}^{-1/2}
 \big|_{\ran H_{r,z,Q}}\right],
\end{equation}
where $H_{r,z,Q}$ and $E_{r,z,Q}$ are the block restrictions and the inverse
acts only on the positive image of $H_{r,z,Q}$. In double precision,
the minimum over all retained nonzero blocks was
\begin{center}
\begin{tabular}{rl}
\toprule
$Q$ & $\min_{r,z}e_{r,z,Q}$\\
\midrule
$10^3$ & $-7.42874\times10^{-4}$\\
$10^5$ & $-6.94088\times10^{-6}$\\
$10^7$ & $-3.58164\times10^{-7}$\\
$10^9$ & $-3.58164\times10^{-7}$\\
Plane & $-3.58164\times10^{-7}$\\
\bottomrule
\end{tabular}
\end{center}
Kernel and descendant-metric residuals were below $1.1\times10^{-15}$ and
$1.8\times10^{-15}$, respectively. The frozen head certificate therefore
does not numerically fit the chosen allowance $2\times10^{-6}$ at $Q=10^5$;
it does at the two larger sampled fluxes. The original conservative
threshold in Appendix~\ref{app:explicit-thresholds} follows from uniform
analytic estimates and exact rational comparisons.

The verifier requires the shifted matrix to be positive semidefinite,
to annihilate the Hamiltonian kernel, and to have the same rank as the
Hamiltonian. The frozen rows pass all three checks.

This candidate proof is accompanied by the finite data and a replayable
checker. It has not been checked by an external referee or
formalized in a proof assistant. The explicit sufficient threshold
$Q_0=3088$ follows from the interval certificates and exact budgets in
Appendix~\ref{app:optimized-thresholds}. No optimality is claimed for this
threshold, which does not locate an onset of the physical gap. At Laughlin flux this sufficient
range starts at $N=1545$. The result concerns the complete projected LLL
interaction on the round sphere and its planar limit.

\paragraph{Source.}
The Fock-space comparison is adapted from OpenAI~\cite{OpenAI2026},
which treats the fermionic $V_1$ interaction. The bosonic coherent-error
bounds and frozen rows used here are given above.

\appendix
\section{The original conservative threshold}\label{app:explicit-thresholds}
This appendix gives the original conservative bound $10^{50}$;
Appendix~\ref{app:optimized-thresholds} gives the current threshold $3088$.
The original argument assigns an explicit sufficient threshold to each
large-$Q$ step in the gap inequality.
 All comparisons are made on fixed few-body spaces before positive
lifting, so their thresholds are uniform in the eventual particle number.
The three thresholds are
\begin{center}
\begin{tabular}{lr}
\toprule
Ingredient & Sufficient integer flux\\
\midrule
Negative three-body tail with allowance $1/2000$ & $Q\ge10^7$\\
Finite coherent comparison bounds $35,2100$ & $Q\ge10^{50}$\\
Frozen head certificate with allowance $2\times10^{-6}$ & $Q\ge10^{50}$\\
\bottomrule
\end{tabular}
\end{center}
Their common sufficient threshold is therefore $Q_0=10^{50}$. This
value comes from the loose norm estimates below. The command
\begin{quote}\ttfamily python -S verify\_explicit\_threshold.py\end{quote}
replays the frozen shifted certificate, the finite planar half-gap, and all
scalar budgets using Python's standard library. The analytic
identities and norm estimates in this appendix complete the argument.
\subsection{An explicit threshold for the negative three-body tail}
\label{app:explicit-tail}

The following estimate replaces the existential threshold of
Theorem~\ref{thm:sphere} at the values needed for the gap budget.

\begin{lemma}\label{lem:explicit-tail}
For every integer $Q\ge10^7$, with the full spherical pair Hamiltonian
normalized as in Section~1,
\begin{equation}\label{eq:explicit-tail-result}
 T_{9,Q}\le\left(c_9+\frac1{2000}\right)(S_{4,Q}+H_Q),
 \qquad c_9=30\left(\frac59\right)^9
             =\frac{19531250}{129140163}.
\end{equation}
\end{lemma}
\begin{proof}
Put
\begin{equation}\label{eq:explicit-tail-parameters}
 a=\frac2{27},\qquad M=\lfloor aQ\rfloor,\qquad
 r=\frac9{10},\qquad t=\frac9{29},\qquad \rho=\frac{29}{36}.
\end{equation}
On the auxiliary space $V_Q\otimes U_Q$, write $\Pi_z$ for the
orthogonal spin-$(3Q/2-z)$ projection, $0\le z\le Q$. Recall the positive
coherent operator
\begin{equation}
 A_{M,Q}=(2Q+1)(Q+1)\int d\Omega_\alpha d\Omega_\beta\,
 s(\alpha,\beta)^M
 |\alpha_{\rm pair}\rangle\langle\alpha_{\rm pair}|
 \otimes|\beta\rangle\langle\beta|
 =\sum_{z=0}^Q\tau_{z,Q}\Pi_z.
\end{equation}
Here each sphere measure has total mass one, the coherent vectors are
normalized, and $s(\alpha,\beta)=\cos^2(\theta_{\alpha\beta}/2)$.
The exact eigenvalue formula is
\begin{align}
 \tau_{z,Q}&=\frac{2Q+1}{3Q-2z+1}
       \sum_{j=z}^Q b_{j,Q}\gamma_{z,j,Q},\label{eq:explicit-tail-tau}\\
 b_{j,Q}&=\frac{Q+1}{M+Q+1}\frac{(Q)_j}{(M+Q)_j},\notag\\
 \gamma_{z,j,Q}&=\frac{(2Q)_z}{(3Q-z+1)_z}
       \binom jz\frac{(Q-z)_{j-z}}{(3Q-2z)_{j-z}}.\notag
\end{align}
The notation $(x)_k=x(x-1)\cdots(x-k+1)$ denotes a descending factorial,
with $(x)_0=1$. Every term in this sum is nonnegative, and the $j=z$
term is positive.

We first bound the finite set $0\le z\le32$ by retaining $j\le100$.
For integers $0\le z\le j\le100$, define
\begin{equation}
 \ell_{z,j}:=\left(\frac23\right)^{z+1}(1+a)^{-j-1}
       \binom jz\left(\frac13\right)^{j-z}
       =2^{z+1}\binom jz t^{j+1}.
\end{equation}
For $Q\ge100$, the $j$th summand in
\eqref{eq:explicit-tail-tau}, including the factor outside the sum, is at
least $\ell_{z,j}P_{z,j,Q}$, where
\begin{equation}\label{eq:explicit-tail-products}
 P_{z,j,Q}=
 \prod_{i=0}^{j-1}\left(1-\frac iQ\right)
 \prod_{i=0}^{z-1}\left(1-\frac i{2Q}\right)
 \prod_{i=z}^{j-1}\left(1-\frac iQ\right).
\end{equation}
To see this factor by factor, use
\begin{equation}
 \frac{2Q+1}{3Q-2z+1}\ge\frac23,\qquad
 \frac{Q+1}{M+Q+1}\ge\frac1{1+a},\qquad
 \frac{Q-i}{M+Q-i}\ge\frac{1-i/Q}{1+a}.
\end{equation}
For $0\le i<z$ and $z\ge1$,
$(2Q-i)/(3Q-z+1-i)\ge(2/3)(1-i/(2Q))$;
for $z=0$ this product is empty. For $0\le i<j-z$,
$(Q-z-i)/(3Q-2z-i)\ge(1/3)(1-(z+i)/Q)$.
All denominators are positive in the stated range. These inequalities
prove~\eqref{eq:explicit-tail-products}.

Every factor in that product belongs to $[0,1]$. The elementary inequality
$\prod_k(1-u_k)\ge1-\sum_k u_k$ for $0\le u_k\le1$, proved by induction,
therefore gives
\begin{equation}
 P_{z,j,Q}\ge
 1-\frac{j(j-1)-z(z-1)/4}{Q}\ge1-\frac{9900}{Q}.
\end{equation}
Consequently, for $0\le z\le32$ and $Q\ge100$,
\begin{equation}\label{eq:explicit-tail-truncated}
 \tau_{z,Q}\ge\left(1-\frac{9900}{Q}\right)
                        \sum_{j=z}^{100}\ell_{z,j}.
\end{equation}

The binomial-series identity, with $0<t<1$, evaluates the full positive sum:
\begin{equation}
 \sum_{j=z}^{\infty}\ell_{z,j}
 =\left(\frac{2t}{1-t}\right)^{z+1}=r^{z+1}.
\end{equation}
Since $\binom jz\le2^j$ and $2t<1$, its omitted part has the uniform
relative bound
\begin{align}
 r^{-z-1}\sum_{j=101}^{\infty}\ell_{z,j}
 &\le\left(\frac2r\right)^{z+1}
       \frac{t}{1-2t}(2t)^{101}\notag\\
 &\le\left(\frac{20}{9}\right)^{33}
       \frac9{11}\left(\frac{18}{29}\right)^{101}
 \le\frac1{1000},\qquad 0\le z\le32.
 \label{eq:explicit-tail-series-check}
\end{align}
The checker specified below verifies the last rational inequality exactly.
For every $Q\ge10^7$, $9900/Q\le1/1000$. Thus
\eqref{eq:explicit-tail-truncated}--\eqref{eq:explicit-tail-series-check}
imply
\begin{equation}\label{eq:explicit-tail-low-z}
 \tau_{z,Q}\ge\left(\frac{999}{1000}\right)^2r^{z+1}
             \ge\frac{499}{500}r^{z+1},\qquad 0\le z\le32.
\end{equation}
Since $(Q-i)/(2Q-i)\le1/2$ for $0\le i<z\le Q$,
$f_{z,Q}=(Q)_z/(2Q)_z\le2^{-z}$. Define
\begin{equation}
 A_*:=\frac2r\left(\frac1{2r}\right)^9
     =\frac{20}{9}\left(\frac59\right)^9=ac_9.
\end{equation}
Since $1/(2r)=5/9<1$, \eqref{eq:explicit-tail-low-z} gives
\begin{equation}\label{eq:explicit-tail-ratio-low}
 \frac{2f_{z,Q}}{\tau_{z,Q}}\le\frac{500}{499}A_*,
 \qquad 9\le z\le32.
\end{equation}

The remaining range is controlled without any expansion in $z/Q$.
Keeping only the $j=z$ term in~\eqref{eq:explicit-tail-tau} gives
\begin{align}
 \frac{2f_{z,Q}}{\tau_{z,Q}}
 &\le2\frac{3Q-2z+1}{2Q+1}\frac{M+Q+1}{Q+1}
 \prod_{i=0}^{z-1}
 \frac{(M+Q-i)(3Q-z+1-i)}{(2Q-i)^2}\notag\\
 &\le3(1+a)\rho^z.\label{eq:explicit-tail-coarse}
\end{align}
For the last inequality the prefactor is at most $3(1+a)$.
With $x=i/Q\in[0,1)$, each product factor is bounded by
\begin{equation}
 \frac{(1+a-x)(3-2x)}{(2-x)^2}\le\frac{3(1+a)}4=\rho.
\end{equation}
After multiplication by the positive denominator, the difference is
$x[(8-4a)-(5-3a)x]/4\ge0$. The second finite rational comparison is
\begin{equation}\label{eq:explicit-tail-coarse-check}
 3\frac{29}{27}\left(\frac{29}{36}\right)^{33}\le A_*.
\end{equation}
Since $0<\rho<1$, this proves $2f_{z,Q}/\tau_{z,Q}\le A_*$ for
every $33\le z\le Q$. In conjunction with
\eqref{eq:explicit-tail-ratio-low}, it gives the whole-space inequality
\begin{equation}
 C_{9,Q}=2\sum_{\substack{9\le z\le Q\\z\text{ odd}}}
                       f_{z,Q}\Pi_z
 \le\frac{500}{499}A_* A_{M,Q},\qquad Q\ge10^7.
\end{equation}

Apply the positive map $X\mapsto\mathcal L_3(W_3XW_3^\dagger)$.
The exact coherent marginal and pair frame give
\begin{equation}
 \mathcal L_3(W_3A_{M,Q}W_3^\dagger)
 \le\frac{Q+1}{M+1}(S_{4,Q}+H_Q).
\end{equation}
Because $M+1>aQ$, for $Q\ge1000$,
\begin{equation}
 \frac{Q+1}{M+1}\le\frac1a\left(1+\frac1Q\right)
                    \le\frac1a\frac{1001}{1000}.
\end{equation}
Using $A_*=ac_9$, the resulting coefficient is at most
$(1001/998)c_9$. Its exact distance below the coefficient in
\eqref{eq:explicit-tail-result} is
\begin{equation}\label{eq:explicit-tail-budget-check}
 c_9+\frac1{2000}-\frac{1001}{998}c_9
 =\frac{1949063779}{42960627558000}>0.
\end{equation}
This proves the claim for every integer $Q\ge10^7$.
\end{proof}

The comparisons in
\eqref{eq:explicit-tail-series-check},
\eqref{eq:explicit-tail-coarse-check}, and
\eqref{eq:explicit-tail-budget-check}, together with the parameter identities
and threshold inequalities, are replayed using integers and
\texttt{fractions.Fraction} by
\path{verify_explicit_tail_threshold.py} with Python's \texttt{-S} flag.
The same argument also permits $Q\ge9900000$; the rounded threshold $10^7$
is used for simplicity.

\subsection{An explicit threshold for the coherent comparison bounds}
\label{app:explicit-coherent-threshold}

We make the coherent comparison bounds of Lemma~\ref{lem:finite-errors}
quantitative, independently of the finite certificate.

\begin{lemma}\label{lem:explicit-coherent-threshold}
For every integer $Q\ge10^{50}$, on every particle-number sector,
\begin{equation}
 \mathcal R_{3,8,Q}\le35(S_{4,Q}+H_Q),\qquad
 \mathcal R_{4,16,Q}\le2100(2S_{4,Q}+H_Q).
\end{equation}
\end{lemma}
\begin{proof}
Put $r\in\{3,4\}$ and $h=r-2$. The auxiliary spin space is
$V_Q\otimes\Sym^hU_Q$: it contains the maximal-spin pair and $h$
spectators, symmetrized only among themselves. Introduce an auxiliary
oscillator space with pair mode $u$ and spectator modes
$v_1,\ldots,v_h$, restricted to spectator symmetry. Its degree-$l$
subspace has normalized occupation labels $(p;j_1,\ldots,j_h)$ with
$p+\sum_sj_s=l$. At $l\le Q$ these labels identify its orthonormal basis
with the spherical auxiliary weight-$l$ basis. All finite-weight matrix
norms below use this unitary identification.

Define the center-of-mass annihilator
$a_\zeta=(\sqrt2a_u+\sum_sa_{v_s})/\sqrt r$. For $z\le16$, let
$F_{z,\infty}$ be an isometric frame from a fixed multiplicity coordinate
space onto the degree-$z$ kernel of $a_\zeta$. Its dimension is one for
$r=3$, and $\lfloor z/2\rfloor+1$ for $r=4$. We need only $z\le8$
when $r=3$. On degree $l$, define the positive diagonal matrix
\begin{equation}\label{eq:explicit-coh-D}
 D_{Q,l}(p;j_1,\ldots,j_h)
 =\sqrt{\frac{\fall{2Q}p}{(2Q)^p}}
       \prod_{s=1}^h\sqrt{\frac{\fall Q{j_s}}{Q^{j_s}}}.
\end{equation}
It preserves spectator symmetry. With $J_\pm$ the total auxiliary spin
raising and lowering operators, the exact matrix identity is
\begin{equation}\label{eq:explicit-coh-raising}
 \frac{J_+}{\sqrt Q}
 =D_{Q,l-1}^{-1}(\sqrt r\,a_\zeta)D_{Q,l}.
\end{equation}
Indeed, the spin raising coefficient for a factor of capacity $cQ$ is
$\sqrt{j(cQ-j+1)}$, where $c=2$ for the pair and $c=1$ for a
spectator; consecutive entries in \eqref{eq:explicit-coh-D} have ratio
$\sqrt{1-(j-1)/(cQ)}$. Thus $D_{Q,z}^{-1}$ maps the oscillator highest
space bijectively onto the spherical highest space. Normalize that map:
\begin{equation}
 V_{z,Q}=D_{Q,z}^{-1}F_{z,\infty},\qquad
 G_{z,Q}=V_{z,Q}^\dagger V_{z,Q},\qquad
 F_{z,Q}=V_{z,Q}G_{z,Q}^{-1/2}.
\end{equation}
Here $G_{z,Q}$ acts in the fixed multiplicity coordinate space.
For $z\le16$, the inequality
$\prod_i(1-x_i)\ge1-\sum_ix_i$, for $x_i\in[0,1]$, gives
\begin{equation}
 (1-128/Q)I\le D_{Q,z}^2\le I.
\end{equation}
The sum of the factors $i/(cQ)$ is at most
$z(z-1)/(2Q)\le120/Q$. For $Q\ge256$ this implies
\begin{equation}\label{eq:explicit-coh-highest-error}
 \|D_{Q,z}^{-1}-I\|\le256/Q,\qquad
 \|G_{z,Q}^{-1/2}-I\|\le256/Q,\qquad
 \|F_{z,Q}-F_{z,\infty}\|\le512/Q.
\end{equation}
For the middle estimate, use
$I\le G_{z,Q}\le(1-128/Q)^{-1}I$ and functional calculus.
The last estimate follows by adding and subtracting
$F_{z,\infty}G_{z,Q}^{-1/2}$.

We next compare normalized descendants through $K=1024$; set $B=4096$.
All retained weights obey $z+n\le1040<B$. Define isometries from the
multiplicity coordinate space to degree $z+n$ by
\begin{equation}
 X_{z,n,Q}=\frac{J_-^nF_{z,Q}}
  {\sqrt{n!\fall{rQ-2z}n}},\qquad
 X_{z,n,\infty}=\frac{(a_\zeta^\dagger)^nF_{z,\infty}}{\sqrt{n!}}.
\end{equation}
The first normalization is the exact spin-$rQ/2-z$ ladder
normalization. The second is isometric because the highest vectors have
center-of-mass occupation zero. On weight $l=z+n$, set
$T_Q=J_-/\sqrt Q$ and $T_\infty=\sqrt r\,a_\zeta^\dagger$.
The respective step maps $X_{z,n+1}=C_{z,n}X_{z,n}$ are
\begin{equation}\label{eq:explicit-coh-step}
 C_{z,n,Q}=\frac{T_Q}{\sqrt{r(n+1)}\sqrt{1-x_{z,n,Q}}},\qquad
 C_{z,n,\infty}=\frac{T_\infty}{\sqrt{r(n+1)}},\qquad
 x_{z,n,Q}=\frac{2z+n}{rQ}.
\end{equation}
In the occupation basis $T_Q$ has nonnegative entries bounded entrywise
by those of $T_\infty$: each one-mode coefficient acquires a factor
$\sqrt{1-i/(cQ)}\le1$. For such matrices,
$|T_Qv|\le T_\infty|v|$ componentwise, which implies the norm bound.
Since $\|T_\infty\|=\sqrt{r(l+1)}$ on degree $l$,
\begin{equation}\label{eq:explicit-coh-step-norm}
 \|C_{z,n,Q}\|
 \le\sqrt{\frac{z+n+1}{n+1}}(1-x_{z,n,Q})^{-1/2}.
\end{equation}
For $Q\ge2^{24}$, the elementary inequality
$1-\sqrt{1-y}\le y$ bounds the sum of the one-mode differences by
$\|T_Q-T_\infty\|\le3B^{3/2}/Q$.
Using $(1-x)^{-1/2}\le2$, $(1-x)^{-1/2}-1\le2x$ for $0\le x\le1/2$,
and $\|C_{z,n,\infty}\|\le\sqrt{17}<5$, gives
\begin{equation}\label{eq:explicit-coh-step-error}
 \|C_{z,n,Q}-C_{z,n,\infty}\|
 \le\frac{6B^{3/2}+10B}{Q}<\frac{2^{22}}Q.
\end{equation}
To control accumulated errors, for $0\le s\le n\le K$ the factors in
\eqref{eq:explicit-coh-step-norm} telescope:
\begin{equation}
 \prod_{t=s}^{n-1}\sqrt{\frac{z+t+1}{t+1}}
 =\sqrt{\frac{\binom{z+n}z}{\binom{z+s}z}}
 \le(z+n)^{z/2}\le2^{88}.
\end{equation}
Empty products equal one, including the case $z=n=0$.
Moreover,
$\sum_{t=s}^{n-1}x_{z,t,Q}\le K(32+K)/(3Q)<2^{19}/Q\le1/32$.
The product inequality used above therefore bounds the product of the
curvature factors $(1-x_{z,t,Q})^{-1/2}$ by $2$.
Every consecutive product of the finite-$Q$ step maps has norm at most
$2^{89}$. Subtracting the two descendant recursions and unrolling them,
\eqref{eq:explicit-coh-highest-error} and
\eqref{eq:explicit-coh-step-error} now give
\begin{equation}\label{eq:explicit-coh-descendant-error}
 \|X_{z,n,Q}-X_{z,n,\infty}\|
 \le\frac{2^{89}(512+K2^{22})}Q<\frac{2^{122}}Q,
 \qquad 0\le n\le K.
\end{equation}

Take $a=1/12$ for $r=3$, and $a=1/16$ for $r=4$, with
$M=\lfloor aQ\rfloor$. The one-spectator north-pole seed $R_Q$ has
diagonal entries
\begin{equation}
 b_{j,Q}=\frac{Q+1}{M+Q+1}\prod_{i=0}^{j-1}\frac{Q-i}{M+Q-i},
 \qquad b_{j,\infty}=(1+a)^{-j-1}.
\end{equation}
Write $M=aQ-e$, $0\le e<1$. For $j\le B$ and $Q\ge2^{24}$,
\begin{equation}
 \left|\frac{Q+1}{M+Q+1}-\frac1{1+a}\right|\le\frac2Q,
 \quad
 \left|\frac{Q-i}{M+Q-i}-\frac1{1+a}\right|
 \le\frac{2(i+1)}Q.
\end{equation}
For the second estimate the difference has numerator $(e-ai)/Q$ and
denominator $(1+a)[1+a-(e+i)/Q]\ge1/2$.
All product factors lie in $[0,1]$, so telescoping the products gives
$|b_{j,Q}-b_{j,\infty}|\le[2+j(j+1)]/Q<2^{26}/Q$.
On degree $l$, define the pair-north seeds
\begin{equation}
 \Sigma_{Q,l}=\left[|0_u\rangle\langle0_u|
          \otimes(R_Q^{\otimes h})_{\rm sym}\right]_l,
 \qquad
 \Sigma_{\infty,l}=\left[|0_u\rangle\langle0_u|
          \otimes(R_\infty^{\otimes h})_{\rm sym}\right]_l.
\end{equation}
Brackets with subscript $l$ mean restriction to degree $l$.
The symmetric spectator basis is orthonormal also for coincident
occupations, so no extra combinatorial factor enters these diagonal
operators. On the retained weights,
\begin{equation}\label{eq:explicit-coh-seed-error}
 \|\Sigma_{Q,l}-\Sigma_{\infty,l}\|\le2^{27}/Q,
 \qquad \|\Sigma_{Q,l}\|,\|\Sigma_{\infty,l}\|\le1.
\end{equation}
For the remaining weights, $Q\ge2/a$ gives $M\ge aQ/2$ and hence
$0\le b_{j,Q}\le(1+a/2)^{-j}$. For both choices of $a$, put $q=32/33$;
then $(1+a/2)^{-1}\le q$, and the limiting seed obeys the same bound.
The pair-north slice at degree $l$ has spectator degree $l$, giving
\begin{equation}\label{eq:explicit-coh-seed-tail}
 \|\Sigma_{Q,l}\|,\|\Sigma_{\infty,l}\|\le q^l.
\end{equation}
If the finite pair-north slice is absent, its seed is zero.

The exact Schur formula for the spin-deficit-$z$ multiplicity matrix is
\begin{equation}
 A_{r,z,Q}=c_{r,z,Q}\sum_{n=0}^{rQ-2z}
       X_{z,n,Q}^\dagger\Sigma_{Q,z+n}X_{z,n,Q},\qquad
 c_{r,z,Q}=\frac{2Q+1}{rQ-2z+1}.
\end{equation}
The oscillator matrix $A_{r,z,\infty}$ is the same sum over $n\ge0$
with the limiting frames and seeds, and $c_{r,z,\infty}=2/r$.
This is the center-of-mass partial trace evaluated by the displacement
twirl in the proof of Lemma~\ref{lem:finite-errors}.
For $Q\ge2^{24}$ and $z\le16$,
$c_{r,z,Q},c_{r,z,\infty}\le1$ and
$|c_{r,z,Q}-c_{r,z,\infty}|\le32/Q$.
Both descendant frames are isometries, so a compressed-seed difference
is bounded by twice the frame error plus the seed error. Thus the
truncated Schur sums differ by at most
\begin{equation}
 \frac{(K+1)(32+2^{27}+2^{123})}Q<\frac{2^{135}}Q.
\end{equation}
The sum of their omitted tails is, by
\eqref{eq:explicit-coh-seed-tail}, at most
$2q^{K+1}/(1-q)=66(32/33)^{1025}$.
The binomial theorem gives $(33/32)^{32}>2$, and therefore
\begin{equation}\label{eq:explicit-coh-block-error}
 \|A_{r,z,Q}-A_{r,z,\infty}\|
 \le\frac{2^{135}}Q+66(32/33)^{1025}
 <\frac{2^{135}}Q+\frac{66}{2^{32}}<10^{-6},
 \qquad Q\ge10^{50}.
\end{equation}
This is uniform in the retained $z$ and all multiplicities.

For $r=3$, the exact limiting eigenvalue is
$[2/(2+3a)]^{z+1}$, whose minimum for $a=1/12$, $z\le8$ is
$t_3=(8/9)^9$. For $r=4$, the eigenvalues are
$(1+a)^{-n_1-1}(1+2a)^{-n_2-1}$, with $n_1+n_2=z$ and $n_1$ even;
their minimum for $a=1/16$, $z\le16$ is
$t_4=(16/17)(8/9)^{17}$. Put $\delta=10^{-6}$.
Equation~\eqref{eq:explicit-coh-block-error}, invariance, and positivity
outside the retained spin blocks imply the whole-space inequalities
\begin{equation}
 \Pi^{\rm aux}_{\le D,Q}\le(t_r-\delta)^{-1}A_{r,Q},
 \qquad (r,D)=(3,8),(4,16).
\end{equation}
Here $A_{3,Q}=A_{M,Q}$ and $A_{4,Q}=A^{(4)}_{M,Q}$ are precisely the
positive kernels of the preceding proofs. Compose with $W_r,W_r^\dagger$
and use the positive lift and exact spectator marginals. Since
$(Q+1)/(M+1)\le a^{-1}(1+1/Q)$ and
$1+1/Q\le10001/10000$ for $Q\ge10^{50}$, the resulting coefficients
are bounded by the exact rational numbers
\begin{equation}
 \frac{12(10001/10000)}{(8/9)^9-10^{-6}}<35,\qquad
 \frac{256(10001/10000)^2}{(16/17)(8/9)^{17}-10^{-6}}<2100.
\end{equation}
All comparisons precede the many-body lift, so the threshold is
independent of particle number. The companion verifier
\path{check_explicit_coherent_threshold.py} checks the final rational
inequalities and the integer bounds used in these constants exactly.
\end{proof}

\subsection{An explicit threshold for transferring the finite certificate}
\label{sec:explicit-head}

The estimates in this subsection give the uniform threshold
$Q_{\rm head}=10^{50}$ for the finite comparison errors.
They use an additional exact calculation on the finite planar
relative spaces retained in the certificate:
\begin{equation}\label{eq:planar-head-half-gap}
 H_{r,\infty}^2-\tfrac12H_{r,\infty}\ge0
 \quad\text{on relative degree }z\le D_r,
 \qquad D_3=8,\quad D_4=16.
\end{equation}
The companion \path{exact_planar_head_gap.py} verifies all 26 rational
quadratic forms and checks that each relative space is invariant under
$H_{r,\infty}$. Since $H_{r,\infty}\ge0$, its nonzero eigenvalues on
these finite spaces are therefore at least $1/2$. 
Throughout this subsection $Q\ge1024$, and all norms are operator norms
in normalized occupation coordinates.

For $n\le4$ and total orbital weight $M\le16$, let $\mathcal W_{n,M}$
be the finite coordinate Hilbert space with orthonormal basis indexed by
partitions $A=(j_1,\ldots,j_n)$ with $\sum_i j_i=M$. These are complete
fixed-weight physical spaces: the same labels identify the spherical and
planar normalized occupation bases for $Q\ge16$. Define the invertible
positive diagonal map $D_{n,M,Q}:\mathcal W_{n,M}\to\mathcal W_{n,M}$ by
\begin{equation}\label{eq:head-diagonal-map}
 D_{n,M,Q}|A\rangle=d_A|A\rangle,
 \qquad d_A=\prod_{j\in A}\sqrt{\prod_{\ell=0}^{j-1}(1-\ell/Q)}.
\end{equation}
It converts normalized spherical coefficients to coefficients of the same
polynomial in the planar monomial convention: after writing
$x=\sqrt Q\,v/u$ and removing $u^Q$, the spherical orbital is
$d_jx^j/\sqrt{j!}$. The map is not unitary. If $J_+$ decreases spherical
orbital weight and $L$ is the planar CM lowering operator, direct action
on the basis gives the exact identity
\begin{equation}\label{eq:head-raising-conjugacy}
 Q^{-1/2}J_+=D_{n,M-1,Q}^{-1}L D_{n,M,Q}.
\end{equation}
For the pair annihilator from $\mathcal W_{n,M}$ to
$\mathcal W_{n-2,M-p}$, the exact Clebsch--Gordan coefficients similarly
give
\begin{equation}\label{eq:head-pair-conjugacy}
 B_{p,Q}=e_{p,Q}^{-1}D_{n-2,M-p,Q}^{-1}
 B_{p,\infty}D_{n,M,Q},\qquad
 e_{p,Q}=\sqrt{\prod_{\ell=0}^{p-1}(1-\ell/(2Q))}.
\end{equation}
The zero-particle diagonal map is the identity.

Fix $r,z$ in the retained ranges, and let
$U:\mathbb C^m\to\mathcal W_{r,z}$ be any isometric frame for $\ker L$,
where $m=\dim\ker L$. The map
$V_{Q,0}=D_{r,z,Q}^{-1}U$ is an invertible frame onto the spherical
highest-weight space, by \eqref{eq:head-raising-conjugacy}. Pullback by
$V_{Q,0}$ is a congruence of quadratic forms. Equation
\eqref{eq:head-pair-conjugacy} shows that their Hamiltonian kernel is the
fixed subspace
\begin{equation}\label{eq:head-fixed-kernel}
 \mathcal N:=\ker(U^\dagger H_{r,\infty}U)
 =\ker(V_{Q,0}^\dagger H_{r,Q}V_{Q,0})\subset\mathbb C^m.
\end{equation}
Indeed both kernels are the simultaneous kernels of all corresponding
pair annihilators. The normal-ordered pair-sandwich structure of
$\mathsf K_r$, and invariance of $\ker H_r$ under rotations and spin
projections, imply that the congruent $E_r$ matrices also annihilate
$\mathcal N$ exactly. Hermiticity then excludes matrix elements between
$\mathcal N$ and $\mathcal N^\perp$.

Since
$\sum_{j\in A}j(j-1)/2\le120$ at weight at most 16, the elementary
product inequality $\prod_i(1-x_i)\ge1-\sum_i x_i$ gives
\begin{equation}\label{eq:head-D-bounds}
 \|D-I\|\le120/Q,\qquad \|D^{-1}\|\le2,\qquad
 \|D^{-1}-I\|\le240/Q,\qquad |e_p^{-1}-1|\le120/Q.
\end{equation}
Each pair interaction is a projector, so
$0\le H_n\le\binom n2I\le6I$ and $\|B_p\|<3$ on either geometry.
Subtracting $B_{p,\infty}$ from \eqref{eq:head-pair-conjugacy} and using
\eqref{eq:head-D-bounds} therefore gives
\begin{equation}\label{eq:head-pair-bound}
 \|B_{p,Q}-B_{p,\infty}\|
 \le\frac{120\cdot2\cdot3+240\cdot3+120\cdot3}{Q}
 <\frac{2000}{Q},\qquad
 \|H_Q-H_\infty\|<\frac{10^6}{Q}.
\end{equation}
The second estimate sums at most 17 pair channels. For the target
$S_3=H_3^2-H_3$, its norm is at most 42 and its difference is at most
$13\cdot10^6/Q$. For $S_4$ on weight $M$, set
$T_p=B_{M-p}^{(2)}B_p^{(4)}$. There are at most 17 such terms,
$\|T_p\|\le9$, and $\|T_{p,Q}-T_{p,\infty}\|\le12000/Q$.
Thus, writing $\mathsf T_3=\mathsf S_3$ and
$\mathsf T_4=(1-c_9-\mu)\mathsf S_4$, where $0<1-c_9-\mu<1$, the
following common bounds hold:
\begin{equation}\label{eq:head-target-bounds}
 \|\mathsf T_{r,Q}\|,\|\mathsf T_{r,\infty}\|<2000,
 \qquad
 \|\mathsf T_{r,Q}-\mathsf T_{r,\infty}\|<10^8/Q.
\end{equation}

The frozen JSON has 138 rows, at most 46 coefficients in each row, and
all its scaled coefficients $y$ have absolute value less than one. The
scaling \eqref{eq:scaledrows}, with $p,j,i\le8$, yields
$|\lambda|\le2$ and
$|\alpha_{pj}|\le\sqrt{2^9\,8!}<5000$. On the output of a pair
annihilator at most two particles remain, and
$\|a_i^\dagger a_j\|\le2$ there. This last bound follows by expressing
the operator as a sum of the single-particle operator
$|i\rangle\langle j|$ in the remaining particle slots. Consequently
\begin{equation}\label{eq:head-row-bounds}
 \|F_Q\|,\|F_\infty\|<2\cdot10^6,\qquad
 \|F_Q-F_\infty\|<10^9/Q.
\end{equation}
For $R_n=\sum_\rho F_\rho^\dagger F_\rho|_n$, these imply
$\|R_n\|<10^{15}$ and $\|R_{n,Q}-R_{n,\infty}\|<10^{18}/Q$.
The normalized $k$-body lift on $n$ particles is the sum over unordered
$k$-particle subsets, hence has norm as a linear map on bounded operators
at most $\binom nk$. It is the same map in matched occupation coordinates.
For the Hermitian inputs needed here, this bound also follows directly
from positivity and $\mathcal L_k^n(I)=\binom nk I$, applied to
$-\|R\|I\le R\le\|R\|I$.
Apply this to the recursive normal-order extraction of $\mathsf K_r$.
The respective multipliers for $r=2,3,4$ are
$1$, $1+3=4$, and $1+6+4\cdot4=23$. For every retained unaveraged
body-weight matrix we therefore obtain
\begin{equation}\label{eq:head-K-bounds}
 \|\mathsf K_{r,M,Q}\|,\|\mathsf K_{r,M,\infty}\|<3\cdot10^{16},
 \qquad
 \|\mathsf K_{r,M,Q}-\mathsf K_{r,M,\infty}\|<3\cdot10^{19}/Q.
\end{equation}
Here body-weight matrices are extended by zero before applying a lift;
the norm of their direct sum is their maximum norm.

For $k=M-z$, define the descendant frames recursively by
\begin{equation}\label{eq:head-descendant-recursion}
 V_{Q,k}=T_{Q,k}V_{Q,k-1},\qquad
 T_{Q,k}:=\frac{J_-/\sqrt Q}
 {\sqrt{k[r-(2z+k-1)/Q]}}.
\end{equation}
Let $U_k$ be the analogous planar frame, starting at $U$ and replacing
$T_{Q,k}$ by $T_{\infty,k}=L^\dagger/\sqrt{rk}$. The SU(2) ladder
identity gives $V_{Q,k}^\dagger V_{Q,k}=V_{Q,0}^\dagger V_{Q,0}$
exactly, so $\|V_{Q,k}\|\le2$, while $\|U_k\|=1$.
On weight $M-1$, $\|L^\dagger\|\le\sqrt{rM}\le8$. Every nonnegative
spherical matrix entry is the corresponding planar entry times
$\sqrt{1-j/Q}$ for $j\le15$, giving
$\|J_-/\sqrt Q-L^\dagger\|\le120/Q$ by entrywise domination.
The inverse denominator in \eqref{eq:head-descendant-recursion} differs
from $1/\sqrt{rk}$ by at most $11/Q$, since
$(2z+k-1)/(rQ)\le31/(3Q)$. The resulting uniform bounds are
\begin{equation}
 \|T_{Q,k}\|\le16,\qquad
 \|T_{Q,k}-T_{\infty,k}\|\le512/Q.
\end{equation}
With $\epsilon_k=\|V_{Q,k}-U_k\|$, we have
$\epsilon_0\le240/Q$ and
$\epsilon_k\le16\epsilon_{k-1}+512/Q$. Since $k\le16$,
\begin{equation}\label{eq:head-descendant-bound}
 \epsilon_k\le\frac{240\cdot16^{16}
       +512(16^{16}-1)/15}{Q}<\frac{10^{22}}Q.
\end{equation}

The exact Schur factor is
$f_Q=(2+1/Q)/(r+(1-2z)/Q)$; its limit is $f_\infty=2/r$.
For the specified ranges, $0<f_Q\le1$ and
$|f_Q-f_\infty|\le16/Q$. If a tilde denotes pullback of a relative
quadratic form by $V_{Q,0}$, the averaged coefficient is exactly
\begin{equation}\label{eq:head-Schur-pullback}
 \widetilde{\mathsf K}_{r,z,Q}
 =f_Q\sum_{M=z}^{D_r}
 V_{Q,M-z}^\dagger\mathsf K_{r,M,Q}V_{Q,M-z}.
\end{equation}
This identity holds for the nonorthonormal highest frame because the
normalized ladder maps preserve its Gram matrix. Combining
\eqref{eq:head-K-bounds} and \eqref{eq:head-descendant-bound}, each
summand differs from its planar counterpart by at most
$[4(3\cdot10^{19})+3(3\cdot10^{16})10^{22}]/Q$.
There are at most 17 summands, and the Schur-factor error contributes
at most $16\cdot17\cdot3\cdot10^{16}/Q$. Hence
\begin{equation}
 \|\widetilde{\mathsf K}_{r,z,Q}
       -\widetilde{\mathsf K}_{r,z,\infty}\|<10^{42}/Q.
\end{equation}
For the highest-frame targets, \eqref{eq:head-D-bounds} gives the bounds
$[4\cdot10^8+3\cdot2000\cdot240]/Q<10^9/Q$ and
$[4\cdot10^6+3\cdot6\cdot240]/Q<5\cdot10^6/Q$ for the target and
Hamiltonian differences, respectively. Thus with $\delta=10^{-6}$ and
$\widetilde F_Q:=\widetilde E_{r,z,Q}
 +2\delta\widetilde H_{r,z,Q}$,
\begin{equation}\label{eq:head-complete-perturbation}
 \|\widetilde F_Q-\widetilde F_\infty\|\le10^{43}/Q.
\end{equation}

The frozen certificate \eqref{eq:exactfinite} and
\eqref{eq:planar-head-half-gap} imply, on the fixed $\mathcal N^\perp$,
\begin{equation}
 \widetilde F_\infty
 =(\widetilde E_\infty+\delta\widetilde H_\infty)
   +\delta\widetilde H_\infty
 \ge\frac\delta2 I=5\cdot10^{-7}I.
\end{equation}
At every integer $Q\ge10^{50}$, the error in
\eqref{eq:head-complete-perturbation} is at most $10^{-7}$.
The remaining margin on $\mathcal N^\perp$ is therefore at least
$4\cdot10^{-7}$, while $\widetilde F_Q$ annihilates $\mathcal N$
exactly. Congruence and rotational invariance give
\begin{equation}\label{eq:explicit-head-result}
 E_{r,Q}\ge-2\cdot10^{-6}H_{r,Q}\Pi^{(r)}_{\le D_r,Q},
 \qquad r=3,4,\quad Q\ge10^{50}.
\end{equation}
The companion \path{explicit_head_threshold_checker.py} checks all scalar
budgets with integers and fractions, validates the frozen row bounds, and
independently recomputes the finite planar half-gap. The planar shifted
certificate itself is replayed by \path{exact_shifted_certificate.py}.


\section{The optimized threshold \texorpdfstring{$Q_0=3088$}{Q0=3088}}\label{app:optimized-thresholds}
The frozen operator rows are unchanged. The optimized argument replaces
the coarse perturbation norms of the previous appendix by direct interval
certification and uses one coherent comparison coefficient for each
retained spin deficit. The real kernel exponents have slopes
$a_3=83/500$ and $a_4=43/500$. These parameters enter only the auxiliary
coherent kernels.

\begin{table}[ht]
\centering\small
\begin{tabular}{rrrrr}
\toprule
Flux interval & $L_3$ & $L_4$ & $A-1/3$ lower bound\\
\midrule
$[3088,3090]$ & $0.00442365$ & $0.00280612$ & $4.2\,10^{-7}$\\
$[3090,3100]$ & $0.00442076$ & $0.00280405$ & $1.3\,10^{-6}$\\
$[3100,3125]$ & $0.004409$   & $0.002802$   & $5.0\,10^{-6}$\\
$[3125,3300]$ & $0.004373$   & $0.002776$   & $6.6\,10^{-7}$\\
$[3300,4000]$ & $0.004139$   & $0.002610$   & $1.93\,10^{-4}$\\
$[4000,\infty)$ & $0.003410$ & $0.002100$   & $4.54\,10^{-4}$\\
\bottomrule
\end{tabular}
\caption{Exact terminating-decimal budgets. For finite upper endpoint
$q_+$ use $g=391/(100q_+)$ in \eqref{eq:finalbudget}; the last row uses
$g=0$. The displayed margin bounds are rounded down and are checked as
rational inequalities. Every row has $B>1/100$.}
\label{tab:optimized-budget}
\end{table}

For each lower endpoint $q_-$ the verifier computes
$\delta_{r,z}=L_r/\overline C_{r,z}(q_-)$ exactly. The coefficient
$\overline C$ is defined below and controls every $Q\ge q_-$.
The head verification likewise covers the entire interval
$0\le1/Q\le1/q_-$, including the plane endpoint. Restricting that already
uniform statement to the row's flux interval permits the sharper tail
budget. The overlapping intervals cover the full flux range.
Six runs each contain all nine three-body and all seventeen four-body
degree blocks, including zero-dimensional ones, for 156 block checks.

Run \path{python -S verify_optimized_threshold.py} from the certificate
directory. It reconstructs the exact budgets, checks all shifts
against the frozen files, replays every interval certificate,
and checks the six final rational budgets and interval coverage.
It also replays the original rational planar certificate and the optimized
tail checker and uses only the Python standard library.

\subsection{An optimized negative-tail bound at every flux}
\label{app:optimized-tail}

Allowing a real coherent weight exponent gives a negative-tail estimate
valid at every flux and a finite-$Q$ improvement.

\begin{lemma}\label{lem:optimized-tail}
Put $c_9=30(5/9)^9$ and $\nu=10^{-10}$. For every integer $Q\ge0$,
\begin{equation}\label{eq:optimized-tail-all-Q}
 T_{9,Q}\le(c_9+\nu)(S_{4,Q}+H_Q).
\end{equation}
The stronger bound
\begin{equation}\label{eq:optimized-tail-improvement}
 T_{9,Q}\le\left(c_9+\nu-\frac dQ\right)(S_{4,Q}+H_Q)
 \qquad(Q\ge Q_*)
\end{equation}
holds for each of the following independent choices:
\begin{equation}\label{eq:optimized-tail-choices}
 (Q_*,d)=(2500,77/20),\qquad(3000,391/100),\qquad(5000,4).
\end{equation}
All these inequalities are independent of total particle number.
\end{lemma}
\begin{proof}
Take $a=2/27$ and $M=aQ$ exactly. The positive coherent auxiliary operator is
\begin{equation}
 A_{M,Q}=(2Q+1)(Q+1)\int d\Omega_\alpha d\Omega_\beta\,
 s(\alpha,\beta)^M
 |\alpha_{\rm pair}\rangle\langle\alpha_{\rm pair}|
 \otimes|\beta\rangle\langle\beta|.
\end{equation}
The measures have total mass one, the vectors are normalized, and
$s(\alpha,\beta)=\cos^2(\theta_{\alpha\beta}/2)$.
Here $M$ is only the exponent of a nonnegative weight, not a spin label;
the integral defines a positive invariant operator for every real $M\ge0$.
The exact beta integral still gives
\begin{equation}
 b_{j,Q}=(Q+1)\binom Qj\int_0^1s^{M+Q-j}(1-s)^j\,ds
 =\frac{Q+1}{M+Q+1}\frac{(Q)_j}{(M+Q)_j}.
\end{equation}
This formula follows by $j$ integrations by parts, valid for
$0\le j\le Q$ and $M\ge0$. The notation $(x)_k$ denotes a descending
factorial, also for real $x$, with $(x)_0=1$.

Writing $A_{M,Q}=\sum_z\tau_{z,Q}\Pi_z$ in the auxiliary pair--single
spin decomposition, the exact eigenvalue formula remains
\begin{align}
 \tau_{z,Q}&=\frac{2Q+1}{3Q-2z+1}
                 \sum_{j=z}^Q b_{j,Q}\gamma_{z,j,Q},\notag\\
 \gamma_{z,j,Q}&=\frac{(2Q)_z}{(3Q-z+1)_z}
                 \binom jz\frac{(Q-z)_{j-z}}{(3Q-2z)_{j-z}}.
 \label{eq:optimized-tail-eigenvalue}
\end{align}
Every term is nonnegative and $\tau_{z,Q}>0$. The exact coherent marginal
and positive lifting give
\begin{equation}
 \mathcal L_3(W_3A_{M,Q}W_3^\dagger)
 \le\frac{Q+1}{M+1}(S_{4,Q}+H_Q).
\end{equation}
Thus it suffices to bound, on every odd $9\le z\le Q$, the scalar cost
\begin{equation}\label{eq:optimized-tail-cost}
 k_{z,Q}:=\frac{Q+1}{M+1}\frac{2f_{z,Q}}{\tau_{z,Q}},
 \qquad f_{z,Q}=\frac{(Q)_z}{(2Q)_z}.
\end{equation}

For $Q>0$, put $x=1/Q$ and
$D_{z,Q}:=2/k_{z,Q}=[(M+1)/(Q+1)]\tau_{z,Q}/f_{z,Q}$.
Dividing each $j=z+n$ term in~\eqref{eq:optimized-tail-eigenvalue} by
$f_{z,Q}$ and cancelling powers of $Q$ gives the exact positive sum
\begin{equation}\label{eq:optimized-tail-D}
 D_{z,Q}=\sum_{n=0}^{Q-z}g_{z,n}(x),
\end{equation}
where
\begin{align}
 g_{z,n}(x)&=\binom{z+n}{z}
 \frac{(2+x)(a+x)}{[3-(2z-1)x](1+a+x)}\notag\\
 &\quad\times\prod_{i=0}^{z-1}
 \frac{(2-ix)^2}{(1+a-ix)[3-(z-1+i)x]}\notag\\
 &\quad\times\prod_{i=0}^{n-1}
 \frac{[1-(z+i)x]^2}{[1+a-(z+i)x][3-(2z+i)x]}.
 \label{eq:optimized-tail-g}
\end{align}
For $Q\ge256$ and $z\le25$, retain the first $33$ terms:
\begin{equation}
 G_z(x):=\sum_{n=0}^{32}g_{z,n}(x)\le D_{z,Q}.
\end{equation}
This truncation is valid because $z+32\le57<Q$. All linear factors in
these rational functions are strictly positive on $[0,1/256]$.

For a positive term
$g(x)=b\prod_i(c_i+d_ix)/\prod_j(e_j+f_jx)$ on $[0,h]$, let the subscript
$-$ or $+$ on a linear factor denote its minimum or maximum endpoint value.
Then
\begin{equation}
 g_-:=\frac{b\prod_i(c_i+d_ix)_-}{\prod_j(e_j+f_jx)_+}
 \le g(x)\le
 g_+:=\frac{b\prod_i(c_i+d_ix)_+}{\prod_j(e_j+f_jx)_-}.
\end{equation}
Since $d/(c+dx)$ has derivative $-d^2/(c+dx)^2\le0$,
\begin{equation}
 \frac{g'(x)}{g(x)}\ge
 L_h(g):=\sum_i\frac{d_i}{c_i+d_ih}-\sum_j\frac{f_j}{e_j}.
\end{equation}
Hence a constant derivative lower bound is
\begin{equation}\label{eq:optimized-tail-interval-rule}
 \beta_h(g):=
 \begin{cases}L_h(g)g_-,&L_h(g)\ge0,\\
               L_h(g)g_+,&L_h(g)<0.
 \end{cases}
 \qquad g'(x)\ge\beta_h(g).
\end{equation}
Repeated factors in~\eqref{eq:optimized-tail-g} are simply listed repeatedly.
All operations and sign decisions in these rules are rational, and the
bounds hold throughout the interval.

Applying these rules to~\eqref{eq:optimized-tail-g}, the supplied exact
checker verifies
\begin{align}
 \sum_{n=0}^{32}\beta_{1/256}(g_{9,n})&>73,
       \label{eq:optimized-tail-derivative}\\
 G_9(0)&\ge\frac2{c_9+10^{-10}},\label{eq:optimized-tail-endpoint}\\
 \sum_{n=0}^{32}(g_{z,n})_-&>20,
 \quad z=11,13,\ldots,25,\quad h=1/256.
 \label{eq:optimized-tail-other}
\end{align}
The endpoint is also checked independently through
\begin{equation}
 G_9(0)=\frac2{c_9}\left(\frac{20}{29}\right)^{10}
         \sum_{n=0}^{32}\binom{n+9}{9}\left(\frac9{29}\right)^n.
\end{equation}
Thus $G_9$ is increasing on the interval, so
$k_{9,Q}\le c_9+\nu$ for every $Q\ge256$. The other retained odd blocks
satisfy $k_{z,Q}<1/10<c_9$ by~\eqref{eq:optimized-tail-other}.

For all remaining spins, the coarse bound obtained by retaining $j=z$ in
\eqref{eq:optimized-tail-eigenvalue} remains valid with real $M=aQ$:
\begin{equation}
 \frac{2f_{z,Q}}{\tau_{z,Q}}\le3(1+a)\rho^z,
 \qquad\rho=\frac{3(1+a)}4=\frac{29}{36}.
\end{equation}
Its factor-by-factor proof uses only $M\le aQ$, not integrality. Since
$(Q+1)/(aQ+1)\le1/a$, every $z\ge27$ obeys
\begin{equation}\label{eq:optimized-tail-high}
 k_{z,Q}\le\frac{3(1+a)}a\left(\frac{29}{36}\right)^{27}
              <\frac{13}{100}<c_9,
\end{equation}
with both comparisons checked exactly.

To exhaust the finite flux remainder, the checker verifies all $2151$
pairs consisting of an integer $9\le Q\le255$ and an odd
$9\le z\le\min(Q,25)$. For each such pair it checks
\begin{equation}
 \sum_{n=0}^{\min(32,Q-z)}g_{z,n}(1/Q)>\frac{100}{7}.
\end{equation}
The omitted terms in~\eqref{eq:optimized-tail-D} are nonnegative, so these
blocks have $k_{z,Q}<7/50<c_9$. For $Q<9$, the negative tail
is empty. These arguments prove~\eqref{eq:optimized-tail-all-Q}.

Finally put $C=c_9+\nu$. The same exact interval rules verify
\begin{equation}\label{eq:optimized-tail-strong-data}
 \sum_{n=0}^{32}\beta_{1/Q_*}(g_{9,n})>b
 \quad\text{for}\quad
 (Q_*,b,d)=(2500,681/2,77/20),\ (3000,345,391/100),\ (5000,353,4).
\end{equation}
For each triple the checker also verifies
\begin{equation}\label{eq:optimized-tail-strong-budget}
 bC^2\ge d\left(2+\frac{bC}{Q_*}\right),
 \qquad C-\frac d{Q_*}>\frac{13}{100}.
\end{equation}
Equations~\eqref{eq:optimized-tail-endpoint} and
\eqref{eq:optimized-tail-strong-data} imply, for every $Q\ge Q_*$,
\begin{equation}
 D_{9,Q}\ge\frac2C+\frac bQ,
 \qquad k_{9,Q}\le\frac{C}{1+bC/(2Q)}\le C-\frac dQ.
\end{equation}
The last inequality follows by cross-multiplication and
\eqref{eq:optimized-tail-strong-budget}. The other odd blocks cost less than
$1/10$ or $13/100$ and therefore fit the same strengthened bound. Positive
lifting now proves~\eqref{eq:optimized-tail-improvement}.
\end{proof}

The tail checker replays all scalar interval certificates and the finite
remainder checks, using only Python's standard library and exact fractions.
Run \path{optimized_tail_threshold_checker.py} with the \texttt{-S} flag.
Its JSON output records the exact derivative
lower bounds, endpoint margins, other retained-block bounds, coarse tail
bound, and finite exhaustion count.

\subsection{A direct quantitative coherent comparison}
\label{app:optimized-coherent}

The following comparison avoids descendant truncation. Throughout this
subsection $r\in\{3,4\}$, $h=r-2$, and the auxiliary space is
$V_Q\otimes\Sym^hU_Q$. Choose any fixed $a>0$ and use the real exponent
$M=aQ$ in the positive coherent kernel. Its positivity and the exact
spectator marginal $(Q+1)/(M+1)$ do not require $M$ to be an integer.

For a retained spin deficit $z$, let
\begin{equation}\label{eq:opt-coh-parameters}
 Z_z=\frac{z(z-1)}2,\qquad k_z=z+r-1,\qquad
 b_z=\frac{a(r/2)^2}{2(1-\theta_z)},\qquad
 c_z=2k_z+b_zk_z(k_z+1),
\end{equation}
where $0<\theta_z<1$ is a proof parameter. Define
\begin{equation}
 t_{3,z}(a)=(1+3a/2)^{-z-1},\qquad
 t_{4,z}(a)=(1+a)^{-1}(1+2a)^{-z-1}.
\end{equation}
These are the exact smallest eigenvalues of the limiting relative
Gaussian kernels in degree $z$, including all multiplicities.

\begin{lemma}\label{lem:opt-coh}
Suppose
\begin{equation}\label{eq:opt-coh-domain}
 Q\ge z,\qquad Q>\max(Z_z,c_z),\qquad
 Q\ge\frac{2(1-\theta_z)}{a\theta_z^2}.
\end{equation}
The spin-deficit-$z$ multiplicity matrix of the coherent kernel obeys
\begin{equation}\label{eq:opt-coh-block-lower}
 A_{r,z,Q}\ge
 \left(1+\frac1{2Q}\right)\left(1+\frac1Q\right)^h
 \left(1-\frac{Z_z}Q\right)\left(1-\frac{c_z}Q\right)t_{r,z}(a)I.
\end{equation}
Consequently its degree-resolved lifted comparison coefficient is
\begin{equation}\label{eq:opt-coh-coefficient}
 C_{r,z}(Q)=
 \frac1{a^ht_{r,z}(a)(1+1/(aQ))^h(1+1/(2Q))
                 (1-Z_z/Q)(1-c_z/Q)}.
\end{equation}
For the three-body lift the corresponding form is $S_{4,Q}+H_Q$;
for the four-body lift it is $2S_{4,Q}+H_Q$.
\end{lemma}
\begin{proof}
Use oscillator coordinates $x=(u,v_1,\ldots,v_h)\in\mathbb C^{h+1}$,
with the spectators symmetrized. Degree-$z$ normalized occupation labels
are identified with the spherical weight-$z$ labels. The normalized
center-of-mass annihilator is
$a_\zeta=(\sqrt2a_u+\sum_sa_{v_s})/\sqrt r$.
Let $F_z$ be an isometric frame for its degree-$z$ kernel. The exact
diagonal matrix
\begin{equation}
 D_{Q,z}(p;j_1,\ldots,j_h)=
 \sqrt{\frac{\fall{2Q}p}{(2Q)^p}}
 \prod_{s=1}^h\sqrt{\frac{\fall Q{j_s}}{Q^{j_s}}},
 \qquad p+\sum_sj_s=z,
\end{equation}
satisfies $J_+/\sqrt Q=D_{Q,z-1}^{-1}(\sqrt r\,a_\zeta)D_{Q,z}$.
Thus $V_z=D_{Q,z}^{-1}F_z$ spans the exact spherical highest space.
With $G_z=V_z^\dagger V_z$, its normalized frame is
$F_{z,Q}=V_zG_z^{-1/2}$. The product inequality
$\prod_i(1-x_i)\ge1-\sum_ix_i$ gives
\begin{equation}\label{eq:opt-coh-gram}
 (1-Z_z/Q)I\le D_{Q,z}^2\le I,
 \qquad I\le G_z\le(1-Z_z/Q)^{-1}I.
\end{equation}
Indeed, the sum of its diagonal product losses is
$[p(p-1)/4+\sum_sj_s(j_s-1)/2]/Q\le Z_z/Q$.

In the stereographic coordinate $\alpha$, the unit-mass spherical measure
is $d^2\alpha/[\pi(1+|\alpha|^2)^2]$, and a spin of capacity $L$ has
coherent coefficients
$(1+|\alpha|^2)^{-L/2}\sqrt{\binom Li}\alpha^i$.
Rescale the pair coordinate to $u=\sqrt{2Q}\alpha$, and the spectators
to $v_s=\sqrt Q\beta_s$. Set
\begin{equation}
 X=|u|^2+\sum_s|v_s|^2,\qquad
 d_s=|v_s-u/\sqrt2|^2,\qquad D_{\rm rel}=\sum_sd_s.
\end{equation}
The relative quadratic form satisfies $D_{\rm rel}\le\kappa X$ with
$\kappa=r/2$: its nonzero eigenvalue is $3/2$ for $r=3$, and its relative
eigenvalues are $1,2$ for $r=4$.

For a multiplicity vector $w$, set $\psi=F_zG_z^{-1/2}w$ and define its
coherent-overlap polynomial
$P(x)=\sum_{\boldsymbol n}\overline{\psi_{\boldsymbol n}}
x^{\boldsymbol n}/\sqrt{\boldsymbol n!}$.
It is homogeneous of degree $z$; the conjugation follows the bra-overlap
convention, and its Gaussian norm equals $\|\psi\|$.
The factor $D_{Q,z}^{-1}$ in the spherical highest vector cancels the
falling-factorial factors of the coherent coefficients exactly. Therefore
\begin{align}
 \langle w,A_{r,z,Q}w\rangle
 &=\mathcal J_Q\int_{\mathbb C^{h+1}}|P(x)|^2W_Q(x)
                      \prod_i\frac{d^2x_i}{\pi},\\
 \mathcal J_Q&=(1+1/(2Q))(1+1/Q)^h,\notag\\
 W_Q(x)&=\prod_{i=0}^h
 \left(1+\frac{|x_i|^2}{\ell_iQ}\right)^{-\ell_iQ-2}
                  \prod_{s=1}^h(1-y_s)^{aQ},\notag\\
 y_s&=\frac{d_s}{Q(1+|u|^2/(2Q))(1+|v_s|^2/Q)},
 \qquad \ell_0=2,\quad\ell_s=1\ (s\ge1).\notag
\end{align}
Here $x_0=u$ and $x_s=v_s$. The identity for $y_s$ follows from
$1-s(\alpha,\beta)=|\alpha-\beta|^2/
[(1+|\alpha|^2)(1+|\beta|^2)]$, so $0\le y_s\le1$.
The normalized spectator symmetrizer gives no additional factor in a
highest-vector overlap.

The corresponding exact planar weight is
$W_\infty=\exp(-X-aD_{\rm rel})$. For $X\le\theta_zQ/\kappa$,
we have $y_s\le d_s/Q\le\theta_z$. The logarithmic-series bound
$-\log(1-y)\le y+y^2/[2(1-\theta_z)]$, together with
$\log(1+t)\le t$ and $\sum_sd_s^2\le\kappa^2X^2$, gives
\begin{equation}\label{eq:opt-coh-global-weight}
 W_Q\ge W_\infty\exp[-(2X+b_zX^2)/Q]
      \ge W_\infty[1-(2X+b_zX^2)/Q].
\end{equation}
For $X>\theta_zQ/\kappa$, the bracket is nonpositive by
\eqref{eq:opt-coh-domain}, since
$b_z(\theta_zQ/\kappa)^2/Q\ge1$.
Positivity of $W_Q$ therefore makes the last bound in
\eqref{eq:opt-coh-global-weight} valid on all of $\mathbb C^{h+1}$.

Normalize the positive measure $|P|^2W_\infty\prod_i d^2x_i/\pi$.
At fixed angular direction $\omega$, writing $x=\sqrt X\omega$ gives
the radial density proportional to
$X^{k_z-1}e^{-\lambda(\omega)X}$, with
$\lambda(\omega)=1+aD_{\rm rel}(\omega)\ge1$.
The exact Gamma-distribution moments therefore give
$\mathbb E[X]\le k_z$ and $\mathbb E[X^2]\le k_z(k_z+1)$,
including after averaging over $\omega$.
Integrating the global weight bound yields
\begin{equation}
 \int|P|^2W_Q\ge(1-c_z/Q)\int|P|^2W_\infty.
\end{equation}
The planar integral is bounded below by
$t_{r,z}(a)\|G_z^{-1/2}w\|^2$.
Equation~\eqref{eq:opt-coh-gram} now proves
\eqref{eq:opt-coh-block-lower}. Inverting this bound, composing with
$W_r,W_r^\dagger$, and applying the positive lift gives the claimed
coefficient, because the exact marginal $\mu_Q=(Q+1)/(aQ+1)$ obeys
\begin{equation}
 \frac{\mu_Q^h}{\mathcal J_Q}
 =\frac{a^{-h}}{(1+1/(aQ))^h(1+1/(2Q))}.
\end{equation}
\end{proof}

For uniform bounds at all $Q\ge Q_0$, define
\begin{equation}\label{eq:opt-coh-uniform}
 \overline C_{r,z}(Q_0)=
 \max\{C_{r,z}(Q_0),a^{-h}t_{r,z}(a)^{-1}\}.
\end{equation}
This bound applies even at small $z$ where $C_{r,z}(Q)$ need not decrease.
Indeed, the logarithm of
$(1+x/a)^h(1+x/2)(1-Z_zx)(1-c_zx)$ is concave for
$0\le x\le1/Q_0$, so its minimum is bounded by its endpoint minimum.
If $Z_z+c_z>h/a+1/2$, logarithmic differentiation gives a strictly
negative derivative, and the simpler bound $C_{r,z}(Q)\le C_{r,z}(Q_0)$
holds for every $Q\ge Q_0$.

In particular, a degree-resolved few-body error satisfying
\begin{equation}
 E_{r,Q}\ge-\sum_{z=0}^D\delta_zH_{r,Q}\Pi^{(r)}_{z,Q},
 \qquad\delta_z\ge0,
\end{equation}
has lifted negative cost at most
$\max_z\delta_z\overline C_{r,z}(Q_0)$ times the corresponding form.
To see this, compare the single auxiliary weighted projector
$\sum_z\delta_z\Pi_z$ with the same positive coherent kernel in every
spin block, and then lift once. The value of $a$ must be common to all
blocks of that kernel; $\theta_z$ may vary since it is only a proof
parameter.

With $a=1/12$ for $r=3$ and $a=1/16$ for $r=4$,
the following thresholds are obtained by exact rational evaluation of
\eqref{eq:opt-coh-coefficient} and its monotonicity:
\begin{center}
\begin{tabular}{ccccc}
\toprule
$r$ & cutoff $D$ & comparison constant & sufficient $Q$ & $\theta$\\
\midrule
3&8&35&4505&$1/14$\\
3&8&36&1254&$1/7$\\
3&8&38&544&$1/5$\\
4&16&2100&4330&$1/12$\\
4&16&2200&2107&$1/8$\\
\bottomrule
\end{tabular}
\end{center}
The threshold in each row is uniform over the retained degrees, since
$Z_z,c_z$ increase with $z$ and $t_{r,z}(a)$ decreases.
The constants $35,2100$ therefore hold simultaneously for
$Q\ge4505$, independently of particle number.
The exact scalar checks are supplied in
\path{optimized_coherent_check.py}; general rational parameters and
degree-resolved bounds are implemented in
\path{optimized_coherent_coefficients.py}.

\subsection{A parameter-interval certificate for the finite head}
\label{sec:optimized-head-interval}

We certify degree-dependent errors on the entire flux interval
$Q\ge q_0$. The parameter is the real scalar $x=Q^{-1}\in[0,X]$, where
$X=q_0^{-1}$; $x=0$ denotes the planar limit.
For each retained pair $(r,z)$ and specified rational $\delta_{r,z}\ge0$,
the verifier proves
\begin{equation}\label{eq:interval-head-claim}
 E_{r,Q}\big|_z+\delta_{r,z}H_{r,Q}\big|_z\ge0
 \qquad\text{for every integer }Q\ge q_0.
\end{equation}
Here restriction is to the highest-weight multiplicity space of deficit
$z$; rotational invariance extends the claim to the whole spin block.
The interval enclosure establishes \eqref{eq:interval-head-claim}
throughout the real parameter interval.

We first express all finite forms in fixed polynomial coordinates. Use
the canonical basis $|A]_{x,Q}=\prod_j(\sqrt{j!}\,a_j^\dagger)^{n_j}|0\rangle$
on the sphere, with $A$ a partition of total weight $M$ into $n$ parts.
The indexed creator $x_j$ from the rational verifier remains distinct from
the scalar $x=Q^{-1}$. The planar Gram entry is
$g_A=\prod_jn_j!(j!)^{n_j}$. Define the scalar polynomials and their
positive square roots by
\begin{equation}\label{eq:interval-scalar-products}
 P_j(x)=\prod_{\ell=0}^{j-1}(1-\ell x),\quad
 T_p(x)=\prod_{\ell=0}^{p-1}(1-\ell x/2),\quad
 d_j=\sqrt{P_j},\quad e_p=\sqrt{T_p},\quad
 d_A=\prod_{j\in A}d_j.
\end{equation}
For $n\le4$ and $M\le16$, let $\mathcal P_{n,M}$ be the coordinate
space of columns indexed by these complete fixed-weight occupation
states. The invertible map to the spherical weight space is
\begin{equation}\label{eq:interval-polynomial-map}
 \Phi_{n,M,Q}:\mathcal P_{n,M}\longrightarrow
       (\Sym^nU_Q)_M,
 \qquad c\longmapsto\sum_Ac_A d_A^{-1}|A]_{x,Q}.
\end{equation}
It is a polynomial-coordinate identification, not a unitary map. Its
Gram matrix is
\begin{equation}\label{eq:interval-metric}
 G_{n,M}(x)=\Phi_{n,M,Q}^\dagger\Phi_{n,M,Q}
          =\operatorname{diag}(g_A/d_A^2).
\end{equation}
All matrices with superscript $\mathrm{pol}$ below are pullbacks of
physical quadratic forms by $\Phi$, already including this metric.

Let $\widehat B_p=\sum_{a=0}^p\partial_a\partial_{p-a}$, with
$\partial_j$ the canonical polynomial derivative. Direct substitution
of the spherical pair coefficients gives the exact identity
\begin{equation}\label{eq:interval-pair-map}
 B_{p,Q}\Phi_{n,M,Q}
 =\frac{\kappa_p}{e_p}\Phi_{n-2,M-p,Q}\widehat B_p,
 \qquad \kappa_p^2=\frac{p!}{2^{p+1}}.
\end{equation}
The polynomial highest-weight space is the fixed rational kernel of
$L=\sum_{j\ge1}j x_{j-1}\partial_j$. The exact raising identity is
\begin{equation}
 Q^{-1/2}J_+\Phi_{n,M,Q}=\Phi_{n,M-1,Q}L.
\end{equation}
If $U$ is a full rational basis of this kernel, then
\begin{equation}\label{eq:interval-common-kernel}
 \ker\big[U^\dagger H^{\mathrm{pol}}_{r,z}(x)U\big]
 =\ker\big[U^\dagger H^{\mathrm{pol}}_{r,z}(0)U\big]
 \quad (0\le x\le X).
\end{equation}
This follows by applying \eqref{eq:interval-pair-map} to every pair
annihilator; all scalar factors and Gram matrices are invertible.
The pair-sandwich structure of the normal-ordered coefficients, together
with invariance of $\ker H_r$ under rotations and spin projections,
implies that the comparison form also annihilates this same kernel exactly.

The verifier chooses a fixed linear complement using principal pivot
coordinates of the exact rational planar Hamiltonian form. Its principal
matrix is positive definite by rational elimination. A further invertible
rational upper triangular change of coordinates approximately normalizes
this Hamiltonian form to the identity. Decimal arithmetic proposes its
entries, which are rounded to exact rationals; the nonzero diagonal is
then checked. Write the resulting canonical coordinate frame as
$U_{\rm a}$. Positivity on its image suffices, because the exact form
annihilates the complementary kernel. A block of rank zero is identically
zero and requires no positive-pivot test.

The interval verifier uses the following direct normal-order construction.
In a fixed channel $t$, label row columns by $0$ and
$a=(p,j)$, put $i=p+j-t$, and define the exact rational covariance
\begin{equation}\label{eq:interval-row-covariance}
 C_{ab}^{(t)}=\sum_\rho y_{\rho,a}y_{\rho,b},
 \qquad h_a(x)=\frac{i!}{e_p d_i d_j}.
\end{equation}
The sum is over the frozen rows in that channel. For $n=3,4$, define the
integer row functional $\ell_{p;l_1,\ldots,l_{n-2}}$ as the vacuum
coefficient of
$\partial_{l_1}\cdots\partial_{l_{n-2}}\widehat B_p$ on the canonical
occupation basis at weight $p+\sum_s l_s$. Integer occupation
multiplicities are included. Within this channel suppress the superscript
$t$ on $C$, and write $b=(p',j')$, $i'=p'+j'-t$. The channel's
unaveraged three-body form is
\begin{align}\label{eq:interval-direct-K3}
 &\sum_a C_{0a}\frac{h_a}{e_t}
   \big(\ell_{t;i}^{\mathsf T}\ell_{p;j}
         +\ell_{p;j}^{\mathsf T}\ell_{t;i}\big)\notag\\
 &\quad+\sum_{a,b:\,i=i'}
    \frac{C_{ab}i!}{e_p e_{p'}d_jd_{j'}}
       \ell_{p;j}^{\mathsf T}\ell_{p';j'}.
\end{align}
Its four-body form is
\begin{equation}\label{eq:interval-direct-K4}
 \sum_{a,b}C_{ab}h_ah_b
       \ell_{p;j,i'}^{\mathsf T}\ell_{p';i,j'}.
\end{equation}
Terms in \eqref{eq:interval-direct-K3} belong to weight $M=p+j$;
terms in \eqref{eq:interval-direct-K4} belong to
$M=p+j+i'=(p+j)+(p'+j')-t$. Summing channels therefore gives the full
forms $\mathsf K^{\mathrm{pol}}_{3,M}$ and
$\mathsf K^{\mathrm{pol}}_{4,M}$, supported at $M\le8$ and $M\le16$.
To verify their coefficients, the physical functional
$a_{l_1}\cdots a_{l_s}B_p$ after pullback is
\begin{equation}
 \frac{\kappa_p\sqrt{l_1!\cdots l_s!}}
 {e_p d_{l_1}\cdots d_{l_s}}\,\ell_{p;l_1,\ldots,l_s}.
\end{equation}
Substitution of $\lambda=y_0/\kappa_t$ and
$\alpha_{p,j}=y_{p,j}\sqrt{i!/(\kappa_p^2j!)}$ into the exact
normal-order identities cancels the square-root factorials and gives
\eqref{eq:interval-direct-K3}--\eqref{eq:interval-direct-K4}.

The remaining targets also have finite coordinate formulas:
\begin{align}\label{eq:interval-target-forms}
 H^{\mathrm{pol}}_{n,M}(x)
 &=\sum_{p=0}^M\frac{\kappa_p^2}{e_p^2}
    \widehat B_p^{\mathsf T}G_{n-2,M-p}(x)\widehat B_p,\notag\\
 \mathsf S^{\mathrm{pol}}_{3,M}
 &=H^{\mathrm{pol}}_{3,M}G_{3,M}^{-1}H^{\mathrm{pol}}_{3,M}
       -H^{\mathrm{pol}}_{3,M},\notag\\
 \mathsf S^{\mathrm{pol}}_{4,M}
 &=\sum_{p=0}^M\frac{\kappa_p^2\kappa_{M-p}^2}
                    {e_p^2 e_{M-p}^2}\,v_p^{\mathsf T}v_p,
 \qquad v_p=\widehat B_{M-p}\widehat B_p.
\end{align}
Here $v_p$ is a row functional to the vacuum; the intermediate diagonal
maps cancel in the two successive pair annihilators. The four-body target
is $(1-c_9-\mu)\mathsf S^{\mathrm{pol}}_{4,M}$.

In polynomial coordinates the rescaled spherical lowering operator is
\begin{equation}\label{eq:interval-descendant-polynomial}
 R(x)=\sum_j(1-jx)x_{j+1}\partial_j,
 \qquad C_k(x)=R(x)^kU_{\rm a},\qquad
 n_k(x)=\prod_{l=1}^k l\,[r-(2z+l-1)x].
\end{equation}
Thus $\Phi C_k/\sqrt{n_k}$ is the normalized descendant of the chosen
highest frame. With $D_3=8,D_4=16$, exact Schur averaging gives the
highest-frame form
\begin{equation}\label{eq:interval-Schur}
 \overline K_{r,z}(x)=\frac{2+x}{r+(1-2z)x}
 \sum_{M=z}^{D_r}
 \frac{C_{M-z}(x)^{\mathsf T}\mathsf K^{\mathrm{pol}}_{r,M}(x)
 C_{M-z}(x)}{n_{M-z}(x)}.
\end{equation}
This formula applies to a nonorthonormal frame because the normalized
ladder maps preserve its Gram matrix. The finite Hermitian matrix tested
by the verifier is
\begin{equation}\label{eq:interval-final-matrix}
 F_{r,z}(x)=U_{\rm a}^{\mathsf T}
  [\mathsf T^{\mathrm{pol}}_{r,z}(x)
     +\delta_{r,z}H^{\mathrm{pol}}_{r,z}(x)]U_{\rm a}
       -\overline K_{r,z}(x),
\end{equation}
where $\mathsf T_3=\mathsf S_3$ and
$\mathsf T_4=(1-c_9-\mu)\mathsf S_4$.

Every scalar interval has integer endpoints divided by
$2^{192}$. Products and reciprocals round lower endpoints down and upper
endpoints up using integer floor and ceiling division. Square roots use
integer square roots of the appropriately scaled endpoints, with an
increment of the upper endpoint unless its square is exact. All
denominator intervals must exclude zero and all square-root intervals
must be nonnegative; otherwise verification stops.

A Taylor interval record of order $m$ encloses the exact representation
\begin{equation}\label{eq:interval-jet-definition}
 f(x)=P_f(x)+x^{m+1}\rho_f(x),\qquad
 P_f(x)=\sum_{k=0}^m c_kx^k,\qquad 0\le x\le X.
\end{equation}
It stores intervals for every fixed coefficient $c_k$ and for all values
of $\rho_f$ on the full interval. The certificates here use $m=3$.
Addition is componentwise. If $\operatorname{tail}_{>m}$ discards
degrees at most $m$ from a polynomial, the product remainder is exactly
\begin{equation}\label{eq:interval-jet-product}
 \rho_{fg}=
 \frac{\operatorname{tail}_{>m}(P_fP_g)}{x^{m+1}}
 +P_f\rho_g+P_g\rho_f+x^{m+1}\rho_f\rho_g.
\end{equation}
The polynomial quotient in the first term is well defined also at zero.
Interval Horner evaluation of this identity encloses the product
remainder on the whole interval.

For an inverse, its Taylor coefficients satisfy
\begin{equation}\label{eq:interval-jet-inverse}
 b_0=c_0^{-1},\qquad
 b_k=-c_0^{-1}\sum_{i=1}^kc_i b_{k-i},\qquad
 \rho_{1/f}=-\frac{\rho_{fP_b}}{f},\quad
 P_b=\sum_{k=0}^mb_kx^k.
\end{equation}
The last identity holds because the polynomial part of $fP_b$ is exactly
one. For the positive square root, the corresponding exact formulas are
\begin{equation}\label{eq:interval-jet-root}
 b_0=\sqrt{c_0},\qquad
 b_k=\frac{c_k-\sum_{i=1}^{k-1}b_ib_{k-i}}{2b_0},\qquad
 \rho_{\sqrt f}=\frac{\rho_f-\rho_{P_b^2}}{\sqrt f+P_b}.
\end{equation}
The remainder identity rationalizes $\sqrt f-P_b$. These identities refer
to the exact coefficients; rounding intervals may enclose a nonzero
width around a canceled coefficient without affecting the algebraic
cancellation. Structural induction using
\eqref{eq:interval-jet-product}--\eqref{eq:interval-jet-root} proves that
the final record encloses every matrix entry throughout $[0,X]$.

After symmetric averaging of the matrix records, write the exact form as
\begin{equation}
 F_{r,z}(x)=\sum_{i=0}^m A_ix^i+x^{m+1}\mathcal E(x).
\end{equation}
All coefficient and remainder matrices are symmetric. If
$\mathcal E_{ij}(x)\in[\ell_{ij},u_{ij}]$, a uniform remainder
bound is
\begin{equation}\label{eq:interval-remainder-norm}
 \gamma=X^{m+1}\max_i\sum_j
              \max\{|\ell_{ij}|,|u_{ij}|\},\qquad
 \|x^{m+1}\mathcal E(x)\|\le\gamma.
\end{equation}
The computed $\gamma$ is rounded upward. This is the maximum absolute
row-sum bound for a symmetric matrix in the chosen finite coordinate
space.

For $s=x/X\in[0,1]$, define the Bernstein control matrices
\begin{equation}\label{eq:interval-Bernstein}
 B_k=\sum_{i=0}^k\frac{\binom ki}{\binom mi}X^i A_i,
 \qquad
 \sum_{i=0}^m A_ix^i
 =\sum_{k=0}^m\binom mk s^k(1-s)^{m-k}B_k.
\end{equation}
The scalar weights are nonnegative and sum to one. The verifier proves
$B_k-\gamma I>0$ by interval $LDL^{\mathsf T}$ elimination, requiring
each pivot interval to have strictly positive lower endpoint. Hence
every actual symmetric matrix enclosed by that interval matrix is
positive definite. If a control matrix is too conservative, repeated
arithmetic averaging performs de Casteljau subdivision into the two
half intervals. The verifier requires both halves to pass, recursively,
using the same valid $\gamma$.

Equations \eqref{eq:interval-remainder-norm} and
\eqref{eq:interval-Bernstein} prove $F_{r,z}(x)\ge0$ for the whole real
interval. The common-kernel identity \eqref{eq:interval-common-kernel}
extends this to the full highest space; invertible coordinate pullback
and rotational invariance give \eqref{eq:interval-head-claim}.

The standard-library script \path{optimized_head_interval.py} takes
\texttt{--qmin}, an exact rational shift JSON through \texttt{--shifts},
and an output path through \texttt{--result}. Its saved records include
the rational frames, every Taylor coefficient and remainder interval,
all Bernstein/LDL pivot intervals, and the exact shifts for the nonzero
blocks. A complete run checks all $9+17=26$ relative blocks.

\begin{thebibliography}{9}

\bibitem{Laughlin1983}
R.~B.~Laughlin,
\emph{Anomalous Quantum Hall Effect: An Incompressible Quantum Fluid with
Fractionally Charged Excitations},
Phys. Rev. Lett. \textbf{50}, 1395--1398 (1983).
\href{https://doi.org/10.1103/PhysRevLett.50.1395}{doi:10.1103/PhysRevLett.50.1395}.

\bibitem{Haldane1983}
F.~D.~M.~Haldane,
\emph{Fractional Quantization of the Hall Effect: A Hierarchy of
Incompressible Quantum Fluid States},
Phys. Rev. Lett. \textbf{51}, 605--608 (1983).
\href{https://doi.org/10.1103/PhysRevLett.51.605}{doi:10.1103/PhysRevLett.51.605}.

\bibitem{OrtizEtAl2013}
G.~Ortiz, Z.~Nussinov, J.~Dukelsky, and A.~Seidel,
\emph{Repulsive interactions in quantum Hall systems as a pairing problem},
Phys. Rev. B \textbf{88}, 165303 (2013).
\href{https://doi.org/10.1103/PhysRevB.88.165303}{doi:10.1103/PhysRevB.88.165303}.

\bibitem{ChenSeidel2015}
L.~Chen and A.~Seidel,
\emph{Algebraic approach to the study of zero modes of Haldane pseudopotentials},
Phys. Rev. B \textbf{91}, 085103 (2015).
\href{https://doi.org/10.1103/PhysRevB.91.085103}{doi:10.1103/PhysRevB.91.085103}.

\bibitem{Rougerie2019}
N.~Rougerie,
\emph{On the Laughlin function and its perturbations},
S\'eminaire Laurent Schwartz---EDP et applications (2018--2019),
Exp. II, 1--17.
\href{https://doi.org/10.5802/slsedp.131}{doi:10.5802/slsedp.131};
\href{https://arxiv.org/abs/1906.11656}{arXiv:1906.11656}.

\bibitem{WarzelYoung2023}
S.~Warzel and A.~Young,
\emph{A Bulk Spectral Gap in the Presence of Edge States for a Truncated
Pseudopotential},
Ann. Henri Poincar\'e \textbf{24}, 133--178 (2023).
\href{https://doi.org/10.1007/s00023-022-01210-z}{doi:10.1007/s00023-022-01210-z}.

\bibitem{OpenAI2026}
OpenAI,
\emph{A Fock-space inequality and the Laughlin spectral gap},
preprint, 24 September 2026.
\href{https://github.com/openai/math/blob/main/preprints/A-Fock-space-inequality-and-the-Laughlin-spectral-gap-September-24-2026/A-Fock-space-inequality-and-the-Laughlin-spectral-gap-September-24-2026.pdf}{Official preprint and source repository}.

\end{thebibliography}

\end{document}

